% =====================================================================
% 第六讲 (Exposé VI) 纤维范畴与下降
% (Catégories fibrées et descente)
% A. Grothendieck
% 原书页码: pp. 116-152
% =====================================================================
\chapter{纤维范畴与下降}

% 与原书一致: 0. Introduction 编号为 0
\setcounter{section}{-1}

\section{引言}

与前一讲引言所宣布的相反，事实证明，不事先在一般范畴中足够仔细地发展下降的语言，就不可能对预概形范畴进行下降，即使在特殊情形下也不行.

"下降"的概念为对象的"粘合"(gluing) 乃至范畴的"粘合"的所有过程提供了普遍框架. 粘合最经典的情形与拓扑空间 $X$ 及用开集 $X_i$ 覆盖 $X$ 有关；若对每个 $i$ 在 $X_i$ 之上给一个纤维空间 (fibre space) $E_i$，并对每个指标对 $(i, j)$ 给一个同构 $f_{ji} \colon E_i|X_{ij} \to E_j|X_{ij}$ (其中 $X_{ij} = X_i \cap X_j$)，满足众所周知的传递条件 (简写为 $f_{kj} f_{ji} = f_{ki}$)，就知道在 $X$ 上存在纤维空间 $E$，它在唯一同构意义下由下列条件确定：存在同构 $f_i \colon E|X_i \xrightarrow{\sim} E_i$，满足关系 $f_{ji} = f_j f_i^{-1}$ (用通常的滥用记号). 设 $X'$ 是诸 $X_i$ 的和空间，它是 $X$ 上的纤维空间 (i.e. 配有连续映射 $X' \to X$). 可以把诸 $E_i$ 的给出更简洁地解释为 $X'$ 上的纤维空间 $E'$，把诸 $f_{ji}$ 解释为 $E'$ 在 $X'' = X' \times_X X'$ 上的两个逆像 $E''_1$ 和 $E''_2$ (按两个典范投影) 之间的同构，于是粘合条件可以写成三重纤维积 $X''' = X' \times_X X' \times_X X'$ 上纤维空间 $E'''_1$ 和 $E'''_3$ 之间的同构的恒等式 (其中 $E'''_i$ 表示 $E'$ 按第 $i$ 个典范投影在 $X'''$ 上的逆像). 从 $E'$ 和 $f$ 构造 $E$ 是"下降"过程的典型情形. 此外，从 $X$ 上的纤维空间 $E$ 出发，若存在 $X$ 的开覆盖 $(X_i)$ 使得诸 $E|X_i$ 同构于 $F \times X_i$，或等价地使 $E$ 在 $X' = \coprod_i X_i$ 上的逆像 $E'$ 同构于 $X' \times F$，则说 $X$ "局部平凡" (locally trivial)，纤维为 $F$.

这样，对象的"粘合"概念与性质的"局部化"概念，都与"基变换" $X' \to X$ 的某类研究有关. 在代数几何 (algebraic geometry) 中，还必须考虑许多其他类型的基变换，特别是忠实平坦态射 $X' \to X$，它们应被看作关于 $X$ 之上"预概形或其他对象"的"局部化"过程. 这类局部化与简单的拓扑局部化 (它不过是其特例) 使用得同样多. 在解析几何 (analytic geometry) 中 (程度较轻) 也是如此.

大多数证明归结为验证，因而省略或仅略述；在需要时我们说明证明中出现的较不显然的图式.

\section{宇宙、范畴、范畴等价}

为避免某些逻辑困难，这里采用"宇宙" (universe) 的概念，它是足够"大"的集合，使通常的集合论运算不越出它；一个"宇宙公理"保证每个对象都位于某个宇宙中. 细节参见 C. Chevalley 与报告人合著的一本正在准备的书\footnote{参见讲 VIII.}. 这样，记号 Ens 不表示所有集合的范畴 (这个观念无意义)，而表示位于给定宇宙中的集合的范畴 (我们在记号中不指明它). 同样，Cat 表示位于该宇宙中的范畴的范畴，Cat 的对象 $X$ 到对象 $Y$ 的"态射"按定义是 $X$ 到 $Y$ 的函子.

若 $\mathcal{C}$ 是范畴，用 $\mathrm{Ob}(\mathcal{C})$ 记 $\mathcal{C}$ 的对象集，用 Fl(C) 记 $\mathcal{C}$ 的箭集 (或 $\mathcal{C}$ 的态射集). 因此写 $X \in \mathrm{Ob}(\mathcal{C})$ 以避免通常的滥用记号 $X \in \mathcal{C}$. 若 $\mathcal{C}$ 和 $\mathcal{C}'$ 是两个范畴，$\mathcal{C}$ 到 $\mathcal{C}'$ 的函子总指通常所说的 $\mathcal{C}$ 到 $\mathcal{C}'$ 的共变函子；它的给出蕴含到达范畴和出发范畴 $\mathcal{C}$，$\mathcal{C}'$ 的给出. $\mathcal{C}$ 到 $\mathcal{C}'$ 的函子构成一个集合，记为 Hom$(\mathcal{C}, \mathcal{C}')$，它是一个记为 Hom$(\mathcal{C}, \mathcal{C}')$ 的范畴的对象集. 按定义，$\mathcal{C}$ 到 $\mathcal{C}'$ 的反变函子是 $\mathcal{C}$ 的对偶范畴 $\mathcal{C}^{\circ}$ 到 $\mathcal{C}'$ 的函子.

我们承认函子 $F \colon \mathcal{S} \to \mathcal{C}$ 的投射极限和归纳极限的概念，特别是这一概念最常见的特殊情形：笛卡儿积与纤维积，直和与余直和 (amalgamated sum) 的对偶概念，以及这些运算的通常形式性质.

例如，在刚才引入的范畴 Cat 中，投射极限 (相对于位于所选定宇宙中的范畴 $\mathcal{I}$) 存在；诸 $\mathcal{C}_i$ 的投射极限范畴 $\mathcal{C}$ 的对象集 (相应地：箭集) 通过取诸范畴 $\mathcal{C}_i$ 的对象集 (相应地：箭集) 的投射极限得到. 最熟知的情形是范畴族的产品. 我们将在下文不断使用两个范畴在第三个之上的纤维积.

关于范畴和函子的一切，在等待已经提到的正在准备的书的同时，参见 \cite{exp06-1} (即使在所论及的一般理论方面也必然很不完备).

借此机会明确范畴等价的概念，它在 \cite{exp06-1} 中阐述得不能令人满意. 函子 $F \colon \mathcal{C} \to \mathcal{C}'$ 称为忠实 (faithful) (相应地：全忠实 (fully faithful))，若对 $\mathcal{C}$ 的每对对象 $S$，$T$，映射 $u \mapsto F(u)$ (从 $\mathrm{Hom}(S, T)$ 到 $\mathrm{Hom}(F(S), F(T))$) 单射 (相应地：双射). 称 $F$ 是范畴等价\index[term]{fanchoudengjia@范畴等价}，若 $F$ 全忠实，且此外 $\mathcal{C}'$ 的每个对象 $S'$ 同构于形如 $F(S)$ 的对象. 证明这等价于：存在 $F$ 的拟逆 (quasi-inverse) 函子 $G \colon \mathcal{C}' \to \mathcal{C}$，i.e. 使 $GF$ 同构于 $\mathrm{id}_{\mathcal{C}}$. 若如此，则 $G$ 的给出和同构 $\varphi \colon GF \to \mathrm{id}_{\mathcal{C}'}$ 等价于：对每个 $S' \in \mathrm{Ob}(\mathcal{C}')$，给出由 $\mathcal{C}$ 的对象 $S$ 和同构 $u \colon F(S) \to S'$ 组成的对 $(S, u)$. (用这些记号，存在唯一函子 $\mathcal{C}' \to \mathcal{C}$，它以所给映射 $S' \mapsto G(S')$ 为对象映射，且使映射 $S' \mapsto \varphi(S')$ 成为函子同态 $FG \to \mathrm{id}_{\mathcal{C}'}$.) 最后，若 $G$ 是 $F$ 的拟逆函子，且选择同构 $\varphi \colon FG \xrightarrow{\sim} \mathrm{id}_{\mathcal{C}}$ 和 $\psi \colon GF \stackrel{\sim}{\longrightarrow} \mathrm{id}_{\mathcal{C}}$，则 \cite[I.1.2]{exp06-1} 中所陈述的关于 $\varphi$，$\psi$ 的两个相容性条件实际上彼此等价，且对任何选定的同构 $\varphi$，存在唯一同构 $\psi$ 使所述条件满足.

\section{一个范畴之上的范畴}

设 $\mathcal{E}$ 是 Univ 中的范畴，它是 Cat 的对象，可以考虑 $\mathcal{E}$ 之上的 Cat 对象的范畴 $\mathbf{Cat}_{/\mathcal{E}}$\index[notations]{$\mathbf{Cat}_{/\mathcal{E}}$}. 这个范畴的对象是函子
\[
p \colon \mathcal{F} \longrightarrow \mathcal{E}.
\]
也说配备这样函子的范畴 $\mathcal{F}$ 是 $\mathcal{E}$ 之上的范畴\index[term]{fanchouzhishangdefanchou@范畴之上的范畴}，或 $\mathcal{E}$-范畴. 把 $\mathcal{E}$ 之上的范畴 $\mathcal{F}$ 到 $\mathcal{E}$ 之上的范畴 $\mathcal{G}$ 的 $\mathcal{E}$-函子称为函子
\[
f \colon \mathcal{F} \longrightarrow \mathcal{G}
\]
使得
\[
q f = p,
\]
其中 $p$ 和 $q$ 分别是 $\mathcal{F}$ 和 $\mathcal{G}$ 的投影函子. 因此 $\mathcal{F}$ 到 $\mathcal{G}$ 的 $\mathcal{E}$-函子 $f$ 的集合与 $\mathbf{Cat}_{/\mathcal{E}}$ 中以 $\mathcal{F}$ 为源、$\mathcal{G}$ 为目标的箭的集合一一对应，尽管没有恒等 (因为像上面那样的 $f$ 的给出不把 $\mathcal{F}$ 和 $\mathcal{G}$ 确定为 $\mathcal{E}$ 之上的范畴)；但当然，如在任何其他范畴 $\mathcal{C}_{/S}$ 中那样，通常滥用语言，把 $\mathcal{E}$-函子 (按上面明确的含义) 等同于范畴 Cat$_{/\mathcal{E}}$ 中的箭.

用
\[
\operatorname{Hom}_{\mathcal{E}}(\mathcal{F}, \mathcal{G})
\]
记 $\mathcal{F}$ 到 $\mathcal{G}$ 的 $\mathcal{E}$-函子的集合. 当然，$\mathcal{E}$-函子的复合是 $\mathcal{E}$-函子 (所考虑的复合按定义对应于 $\mathbf{Cat}_{/\mathcal{E}}$ 中箭的复合).

现在考虑两个 $\mathcal{E}$-函子
\[
f, g \colon \mathcal{F} \longrightarrow \mathcal{G}
\]
和函子同态
\[
u \colon f \longrightarrow g.
\]
称 $u$ 为 $\mathcal{E}$-同态或"$\mathcal{E}$-函子同态"，若对每个 $\xi \in \mathrm{Ob}(\mathcal{F})$，有
\[
q(u(\xi)) = \mathrm{id}_{p(\xi)},
\]
用语言说：令 $S = p(\xi) = qf(\xi) = qg(\xi) \in \mathrm{Ob}(\mathcal{E})$，$\mathcal{G}$ 中的态射
\[
u(\xi) \colon f(\xi) \longrightarrow g(\xi)
\]
是 $\mathrm{id}_S$-态射. (一般地，对 $\mathcal{E}$ 中的每个态射 $\alpha \colon T \to S$ 和 $\mathcal{E}$ 之上的每个范畴 $\mathcal{G}$，若 $q(v) = \alpha$ (其中 $q$ 记投影函子 $\mathcal{G} \to \mathcal{E}$)，则 $\mathcal{G}$ 中的态射 $v$ 称为 $\alpha$-态射.) 若有第三个 $\mathcal{E}$-函子 $h \colon \mathcal{F} \to \mathcal{G}$ 和 $\mathcal{E}$-同态 $v \colon g \to h$，则 $vu$ 也是 $\mathcal{E}$-同态. 这样，$\mathcal{F}$ 到 $\mathcal{G}$ 的 $\mathcal{E}$-函子及它们的 $\mathcal{E}$-同态构成 $\mathcal{F}$ 到 $\mathcal{G}$ 的所有函子的范畴 Hom$(\mathcal{F}, \mathcal{G})$ 的一个子范畴，称它为 $\mathcal{F}$ 到 $\mathcal{G}$ 的 $\mathcal{E}$-函子范畴，记为
\[
\mathbf{Hom}_{\mathcal{E}/-}(\mathcal{F}, \mathcal{G}).
\]
它也是函子对
\[
R, S \colon \operatorname{Hom}(\mathcal{F}, \mathcal{G}) \Longrightarrow \operatorname{Hom}(\mathcal{F}, \mathcal{E})
\]
的核子范畴，其中 $R$ 是由 Hom$(\mathcal{F}, \mathcal{E})$ 的对象 $p$ 定义的常函子，$S$ 是由 $q \colon \mathcal{G} \to \mathcal{E}$ 定义的函子 $f \mapsto q \circ f$.

为了结束这些一般理论，还需定义范畴 Hom$_{\mathcal{E}/-}(\mathcal{F}, \mathcal{G})$ 之间由 $\mathcal{E}$-函子复合给出的自然配对. 换言之，想定义"复合函子"：
\begin{equation}
\mathbf{Hom}_{\mathcal{E}/-}(\mathcal{F}, \mathcal{G}) \times \mathbf{Hom}_{\mathcal{E}/-}(\mathcal{G}, \mathcal{H}) \longrightarrow \mathbf{Hom}_{\mathcal{E}/-}(\mathcal{F}, \mathcal{H}) \tag{i}
\label{eq:VI.i}
\end{equation}
(当 $\mathcal{F}, \mathcal{G}, \mathcal{H}$ 是 $\mathcal{E}$ 之上的三个范畴时)，使该函子对对象给出 $\mathcal{E}$-函子 $f \colon \mathcal{F} \to \mathcal{G}$ 和 $g \colon \mathcal{G} \to \mathcal{H}$ 的复合映射 $(f, g) \mapsto gf$. 为此，回忆定义典范函子：
\begin{equation}
\operatorname{Hom}(\mathcal{F}, \mathcal{G}) \times \operatorname{Hom}(\mathcal{G}, \mathcal{H}) \longrightarrow \operatorname{Hom}(\mathcal{F}, \mathcal{H}) \tag{ii}
\label{eq:VI.ii}
\end{equation}
它对对象正是函子复合映射 $(f, g) \mapsto gf$，并把箭 $(u, v)$ (其中
\[
u \colon f \longrightarrow f', \quad v \colon g \longrightarrow g'
\]
分别是 Hom$(\mathcal{F}, \mathcal{G})$ 和 Hom$(\mathcal{G}, \mathcal{H})$ 中的箭) 变为箭
\[
v * u \colon gf \longrightarrow g'f'
\]
由关系定义：
\[
v * u(\xi) = v(f'(\xi)) \cdot g(u(\xi)) = g'(u(\xi)) \cdot v(f(\xi))
\]
众所周知由此确实得到 $gf$ 到 $g'f'$ 的同态，且 (对变动的 $f, g$ 和 $u, v$) 得到函子 (ii)，i.e. 有
\begin{enumerate}[(I)]
\item $\mathrm{id}_g * \mathrm{id}_f = \mathrm{id}_{gf}$，
\item $(v' * u') \circ (v * u) = (v' \circ v) * (u' \circ u)$.
\end{enumerate}
还回忆典范配对 (ii) 的结合律公式，它一方面由函子复合的结合性 $(hg)f = h(gf)$ 表达，另一方面由公式
\begin{equation}
(w * v) * u = w * (v * u) \tag{III}
\label{eq:VI.III}
\end{equation}
表达 (其中 $u \colon f \to f'$，$v \colon g \to g'$ 如上，且还给出函子 $h, h' \colon \mathcal{H} \to \mathcal{H}$ 的同态 $w \colon h \to h'$). 现在说，当 $\mathcal{F}, \mathcal{G}$ 是 $\mathcal{E}$-范畴时，典范复合函子 (ii) 诱导函子 (i). 由于已知两个 $\mathcal{E}$-函子的复合是 $\mathcal{E}$-函子，这等于说：当 $u \colon f \to f'$ 和 $v \colon g \to g'$ 是 $\mathcal{E}$-函子同态时，$v * u \colon gf \to g'f'$ 也是 $\mathcal{E}$-函子同态. 这确实由定义平凡得出. 由于配对 (i) 由配对 (ii) 诱导，它们满足同样的结合律性质，也用公式 $(hg)f = h(gf)$ 和 $(w * v) * u = w * (v * u)$ 表达 (对 $\mathcal{E}$-函子和 $\mathcal{E}$-函子同态).

为完成公式表 (I)，(II)，(III)，还回忆公式：
\begin{equation}
v * \mathrm{id}_{\mathcal{F}} = v \quad \text{与} \quad \mathrm{id}_{\mathcal{G}} * u = u, \tag{IV}
\label{eq:VI.IV}
\end{equation}
其中为简化，当 $u$ 或 $v$ 分别是 $f$ 或 $g$ 的恒等自同构时，用 $v * f$ 或 $u * g$ 代替 $v * u$.

由配对 (i) 的定义得出，Hom$_{\mathcal{E}/-}(\mathcal{F}, \mathcal{G})$ 是 $\mathcal{F}, \mathcal{G}$ 的函子，从范畴 $\mathbf{Cat}_{/\mathcal{E}}^{\circ} \times \mathbf{Cat}_{/\mathcal{E}}$ 到 Cat. 若确实 $g \colon \mathcal{G} \to \mathcal{G}_1$ 是 $\mathcal{E}$-函子，i.e. Hom$_{\mathcal{E}/-}(\mathcal{G}, \mathcal{G}_1)$ 的对象，则在 (i) 中取 $\mathcal{H} = \mathcal{G}_1$，它对应函子
\[
g_{*} \colon \operatorname{Hom}_{\mathcal{E}/-}(\mathcal{F}, \mathcal{G}) \longrightarrow \operatorname{Hom}_{\mathcal{E}/-}(\mathcal{F}, \mathcal{G}_1).
\]
类似地，对 $\mathcal{E}$-函子 $f \colon \mathcal{F}_1 \to \mathcal{F}$，定义函子
\[
f^{*} \colon \operatorname{Hom}_{\mathcal{E}/-}(\mathcal{F}, \mathcal{G}) \longrightarrow \operatorname{Hom}_{\mathcal{E}/-}(\mathcal{F}_1, \mathcal{G}).
\]
为简洁，这些函子也用记号 $f \mapsto g \circ f$ (相应地：$g \mapsto g \circ f$) 表示 (它们事实上只表示对象集上的相应映射). 由上面指出的结合律性质得出，如此确实得到如所宣布的函子 $\mathbf{Cat}_{/\mathcal{E}}^{\circ} \times \mathbf{Cat}_{/\mathcal{E}} \to \mathbf{Cat}$.

\section{\texorpdfstring{$\mathcal{E}$}{E} 之上范畴中的基变换}

由于在 Cat 中投射极限 (相对于 Univ 的元素的范畴 $\mathcal{I}$) 存在，在 $\mathbf{Cat}_{/\mathcal{E}}$ 中也如此，特别笛卡儿积存在，它解释为 Cat 中的纤维积. 按照一般记号，若 $\mathcal{F}$ 和 $\mathcal{G}$ 是 $\mathcal{E}$ 之上的范畴，用
\[
\mathcal{F} \times_{\mathcal{E}} \mathcal{G}
\]
记它们在 $\mathbf{Cat}_{/\mathcal{E}}$ 中的积，i.e. 它们在 Cat 中 $\mathcal{E}$ 之上的纤维积，作为 $\mathcal{E}$ 之上的范畴. 因此 $\mathcal{F} \times_{\mathcal{E}} \mathcal{G}$\index[notations]{$\mathcal{F} \times_{\mathcal{E}} \mathcal{G}$} 配备两个 $\mathcal{E}$-函子 $\mathrm{pr}_1$ 和 $\mathrm{pr}_2$，它们对 $\mathcal{E}$ 之上的每个范畴 $\mathcal{H}$ 定义双射
\[
\operatorname{Hom}_{\mathcal{E}}(\mathcal{H}, \mathcal{F} \times_{\mathcal{E}} \mathcal{G}) \stackrel{\sim}{\longrightarrow} \operatorname{Hom}_{\mathcal{E}}(\mathcal{H}, \mathcal{F}) \times \operatorname{Hom}_{\mathcal{E}}(\mathcal{H}, \mathcal{G}).
\]
这个双射还来自范畴同构
\[
\mathbf{Hom}_{\mathcal{E}/-}(\mathcal{H}, \mathcal{F} \times_{\mathcal{E}} \mathcal{G}) \stackrel{\sim}{\longrightarrow} \mathbf{Hom}_{\mathcal{E}/-}(\mathcal{H}, \mathcal{F}) \times \mathbf{Hom}_{\mathcal{E}/-}(\mathcal{H}, \mathcal{G})
\]
(取两个成员的对象集)，其中所写的函子是以 $h \mapsto \mathrm{pr}_1 \circ h$ 和 $h \mapsto \mathrm{pr}_2 \circ h$ 为分量的函子. 验证确实得到同构留给读者 (在范畴的投射极限的更一般情形也成立，而不限于纤维积).

回忆 (如第 1 号中所说)
\[
\operatorname{Ob}(\mathcal{F} \times_{\mathcal{E}} \mathcal{G}) = \operatorname{Ob}(\mathcal{F}) \times_{\operatorname{Ob}(\mathcal{E})} \operatorname{Ob}(\mathcal{G})
\]
\[
\operatorname{Fl}(\mathcal{F} \times_{\mathcal{E}} \mathcal{G}) = \operatorname{Fl}(\mathcal{F}) \times_{\operatorname{Fl}(\mathcal{E})} \operatorname{Fl}(\mathcal{G}),
\]
箭的复合逐分量进行.

在下面，考虑函子
\[
\lambda \colon \mathcal{E}' \longrightarrow \mathcal{E}
\]
并对 $\mathcal{E}$ 之上的每个范畴 $\mathcal{F}$，借助 $\mathrm{pr}_2$ 把 $\mathcal{F} \times_{\mathcal{E}} \mathcal{E}'$ 看作 $\mathcal{E}'$ 之上的范畴；换言之，把"纤维积"运算解释为"基变换"运算，函子 $\lambda \colon \mathcal{E}' \to \mathcal{E}$ 称为"基变换函子". 按照众所周知的一般事实，得到函子，称为 $\lambda$ 的基变换函子：
\[
\lambda^{*} \colon \mathbf{Cat}_{\mathcal{E}}\index[notations]{$\lambda^{*}$} \longrightarrow \mathbf{Cat}_{\mathcal{E}'},
\]
("基限制"函子的伴随，后者给 $\mathcal{E}'$ 之上的范畴 $\mathcal{F}'$ (投影函子为 $p'$) 对应由函子 $p = \lambda p'$ 看作 $\mathcal{E}$ 之上范畴的 $\mathcal{F}'$). 如范畴中基变换函子的一般情形所熟知，基变换函子"与投射极限可交换"，特别"把" $\mathcal{E}$ 上的纤维积"变为" $\mathcal{E}'$ 上的纤维积.

设 $\mathcal{F}$ 和 $\mathcal{G}$ 是 $\mathcal{E}$ 之上的两个范畴；定义典范同构：
\begin{equation}
\mathbf{Hom}_{\mathcal{E}'/-}(\mathcal{F}', \mathcal{G}') \xrightarrow{\sim} \mathbf{Hom}_{\mathcal{E}/-}(\mathcal{F} \times_{\mathcal{E}} \mathcal{E}', \mathcal{G}) \tag{i}
\label{eq:VI.i2}
\end{equation}
其中 $\mathcal{F}' = \mathcal{F} \times_{\mathcal{E}} \mathcal{E}'$，$\mathcal{G}' = \mathcal{G} \times_{\mathcal{E}} \mathcal{E}'$.

为此考虑函子
\[
\mathrm{pr}_1 \colon \mathcal{G}' = \mathcal{G} \times_{\mathcal{E}} \mathcal{E}' \longrightarrow \mathcal{G},
\]
并用
\[
F \longmapsto \mathrm{pr}_1 \circ F
\]
定义 (i)，它先验地表示函子
\[
\operatorname{Hom}(\mathcal{F}', \mathcal{G}') \longrightarrow \operatorname{Hom}(\mathcal{F}', \mathcal{G})
\]
只需验证后者对子范畴诱导函子 (i)，且它是同构. (ii) 诱导双射
\[
\operatorname{Hom}_{\mathcal{E}'/-}(\mathcal{F}', \mathcal{G}') \stackrel{\sim}{\longrightarrow} \operatorname{Hom}_{\mathcal{E}/-}(\mathcal{F} \times_{\mathcal{E}} \mathcal{E}', \mathcal{G})
\]
是基变换函子的特征性质. 剩下证明：若 $F$，$G$ 是 $\mathcal{F}' \to \mathcal{G}'$ 的 $\mathcal{E}'$-函子，则映射
\[
u \longmapsto \mathrm{pr}_1 \circ u
\]
诱导双射
\[
\operatorname{Hom}_{\mathcal{E}'}(F, G) \stackrel{\sim}{\longrightarrow} \operatorname{Hom}_{\mathcal{E}}(\mathrm{pr}_1 \circ F, \mathrm{pr}_1 \circ G).
\]
这个事实的验证是直接的，留给读者.

由这个同构 (i) 和前一节末尾，得出
\[
\mathbf{Hom}_{\mathcal{E}'/-}(\mathcal{F} \times_{\mathcal{E}} \mathcal{E}', \mathcal{G} \times_{\mathcal{E}} \mathcal{E}')
\]
可以看作 $\mathcal{E}', \mathcal{F}, \mathcal{G}$ 的函子，从范畴 $\mathbf{Cat}_{/\mathcal{E}}^{\circ} \times \mathbf{Cat}_{/\mathcal{E}}^{\circ} \times \mathbf{Cat}_{/\mathcal{E}}$ 到 Cat，同构于由表达式 Hom$_{\mathcal{E}/-}(\mathcal{F} \times_{\mathcal{E}} \mathcal{E}', \mathcal{G})$ 定义的函子. 特别地，对固定的 $\mathcal{F}, \mathcal{G}$，得到 $\mathcal{E}'$ 的函子 $\mathbf{Hom}_{\mathcal{E}'/-}(\mathcal{F}', \mathcal{G}') = \mathbf{Hom}_{\mathcal{E}/-}(\mathcal{F} \times_{\mathcal{E}} \mathcal{E}', \mathcal{G} \times_{\mathcal{E}} \mathcal{E}')$，特别投影 $\mathcal{E}$-函子 $\lambda \colon \mathcal{E}' \to \mathcal{E}$ 定义态射，i.e. 函子
\[
\lambda_{\mathcal{F},\mathcal{G}}^{*} \colon \operatorname{Hom}_{\mathcal{E}/-}(\mathcal{F}, \mathcal{G}) \longrightarrow \operatorname{Hom}_{\mathcal{E}'/-}(\mathcal{F}', \mathcal{G}')
\]
我们把它明确化. 对两个成员的对象集，它是映射
\[
f \longmapsto f' = f \times_{\mathcal{E}} \mathcal{E}'
\]
它表达 $\mathcal{F} \times_{\mathcal{E}} \mathcal{E}'$ 对 $\mathcal{E}$ 上的对象 $\mathcal{F}$ 的函子性依赖. 另一方面，考虑两个 $\mathcal{E}$-函子
\[
f, g \colon \mathcal{F} \longrightarrow \mathcal{G}
\]
和 $\mathcal{E}$-函子同态
\[
u \colon f \longrightarrow g,
\]
明确相应的 $\mathcal{E}'$-函子同态：
\[
u' \colon f' \longrightarrow g'.
\]
对每个
\[
\xi' = (\xi, S') \in \operatorname{Ob}(\mathcal{F}')
\]
(其中
\[
\xi \in \operatorname{Ob}(\mathcal{F}), \quad S' \in \operatorname{Ob}(\mathcal{E}'), \quad p(\xi) = \lambda(S') = S
\])
态射
\[
u'(\xi') \colon f'(\xi') = (f(\xi), S') \longrightarrow g'(\xi') = (g(\xi), S') \quad \text{在 } \mathcal{G}' \text{ 中}
\]
由公式定义：
\[
u'(\xi') = (u(\xi), \mathrm{id}_{S'})
\]
(这确实是 $\mathcal{G}'$ 中的 $S'$-态射，因为 $q(u(\xi)) = \lambda(\mathrm{id}_{S'}) = \mathrm{id}_S$).

现在考虑任意 $\mathcal{E}$-函子
\[
\lambda' \colon \mathcal{E}'' \longrightarrow \mathcal{E}'
\]
和相应函子
\[
\mathbf{Hom}_{\mathcal{E}'/-}(\mathcal{F} \times_{\mathcal{E}} \mathcal{E}', \mathcal{G} \times_{\mathcal{E}} \mathcal{E}') \longrightarrow \mathbf{Hom}_{\mathcal{E}''/-}(\mathcal{F} \times_{\mathcal{E}} \mathcal{E}'', \mathcal{G} \times_{\mathcal{E}} \mathcal{E}'')
\]
我说这个函子正是用前述方法从 $\mathcal{E}'$ 上的 $\mathcal{F}'$ 和 $\mathcal{G}'$ 出发、把 $\mathcal{E}''$ 看作 $\mathcal{E}'$ 之上的范畴得到的函子，考虑到"基变换传递性"同构：
\[
\mathcal{F}' \times_{\mathcal{E}'} \mathcal{E}'' \stackrel{\sim}{\longrightarrow} \mathcal{F}'' = \mathcal{F} \times_{\mathcal{E}} \mathcal{E}'' \quad \text{与} \quad \mathcal{G}' \times_{\mathcal{E}'} \mathcal{E}'' \stackrel{\sim}{\longrightarrow} \mathcal{G}'' = \mathcal{G} \times_{\mathcal{E}} \mathcal{E}'',
\]
它蕴含典范同构
\[
\mathbf{Hom}_{\mathcal{E}''/-}(\mathcal{F}' \times_{\mathcal{E}'} \mathcal{E}'', \mathcal{G}' \times_{\mathcal{E}'} \mathcal{E}'') \stackrel{\sim}{\longrightarrow} \mathbf{Hom}_{\mathcal{E}''/-}(\mathcal{F} \times_{\mathcal{E}} \mathcal{E}'', \mathcal{G} \times_{\mathcal{E}} \mathcal{E}'')
\]
这个相容性的验证是直接的，留给读者.

刚才定义的函子与前一节定义的配对相容，精确地说，若 $\mathcal{F}, \mathcal{G}, \mathcal{H}$ 是 $\mathcal{E}$ 之上的范畴且令
\[
\mathcal{F}' = \mathcal{F} \times_{\mathcal{E}} \mathcal{E}', \quad \mathcal{G}' = \mathcal{G} \times_{\mathcal{E}} \mathcal{E}', \quad \mathcal{H}' = \mathcal{H} \times_{\mathcal{E}} \mathcal{E}',
\]
则在函子图式中有交换性：
\[
\begin{tikzcd}[row sep=large, column sep=large]
\mathbf{Hom}_{\mathcal{E}/-}(\mathcal{F}, \mathcal{G}) \times \mathbf{Hom}_{\mathcal{E}/-}(\mathcal{G}, \mathcal{H}) \arrow[r] \arrow[d,"{\lambda_{\mathcal{F},\mathcal{G}}^{*} \times \lambda_{\mathcal{G},\mathcal{H}}^{*}}"] & \mathbf{Hom}_{\mathcal{E}/-}(\mathcal{F}, \mathcal{H}) \arrow[d,"{\lambda_{\mathcal{F},\mathcal{H}}^{*}}"] \\
\mathbf{Hom}_{\mathcal{E}'/-}(\mathcal{F}', \mathcal{G}') \times \mathbf{Hom}_{\mathcal{E}'/-}(\mathcal{G}', \mathcal{H}') \arrow[r] & \mathbf{Hom}_{\mathcal{E}'/-}(\mathcal{F}', \mathcal{H}')
\end{tikzcd}
\]
其中水平箭是前一节定义的复合函子. 这个交换性用公式表达：
\[
(gf)' = g'f'
\]
(对 $f \in \mathrm{Hom}_{\mathcal{E}}(\mathcal{F}, \mathcal{G})$，$g \in \mathrm{Hom}_{\mathcal{E}}(\mathcal{G}, \mathcal{H})$，该公式简单表达基变换的函子性)，以及公式
\[
(v * u)' = v' * u'
\]
(当 $u \colon f \to f_1$ 是 Hom$_{\mathcal{E}/-}(\mathcal{F}, \mathcal{G})$ 的箭，$v \colon g \to g_1$ 是 Hom$_{\mathcal{E}/-}(\mathcal{G}, \mathcal{H})$ 的箭时). 这个公式的验证由定义容易得出.

下面，我们主要关心 Hom$_{\mathcal{E}}(\mathcal{F}, \mathcal{G})$ (及其些值得注意的子范畴) 当 $\mathcal{F} = \mathcal{E}$ 的情形，为此引入特殊记号：
\[
\boldsymbol{\Gamma}(\mathcal{G}/\mathcal{E}) = \mathbf{Hom}_{\mathcal{E}}(\mathcal{E}, \mathcal{G}), \quad \Gamma(\mathcal{G}/\mathcal{E}) = \mathrm{Ob}\index[notations]{$\Gamma(\mathcal{G}/\mathcal{E})$}(\boldsymbol{\Gamma}(\mathcal{G}/\mathcal{E})) = \mathrm{Hom}_{\mathcal{E}}(\mathcal{E}, \mathcal{G}).
\]

\begin{remarkstar}
当 $\mathcal{E}$ 是单点范畴，i.e. $\mathrm{Ob}(\mathcal{E})$ 和 $\mathrm{Fl}(\mathcal{E})$ 约化为一个元素，这也就意味着 $\mathcal{E}$ 是 Cat 的终对象，则 $\mathcal{E}$ 之上范畴的给出等价于一个范畴的给出 (因为将有 $\mathcal{F}$ 到 $\mathcal{E}$ 的唯一函子). 更精确地说，$\mathbf{Cat}_{/\mathcal{E}}$ 此时同构于 Cat. 此外，范畴 Hom$_{\mathcal{C}/-}(\mathcal{F}, \mathcal{G})$ 此时正是 $\mathbf{Hom}(\mathcal{F}, \mathcal{G})$. 回忆基本公式
\[
\operatorname{Hom}(\mathcal{H}, \operatorname{Hom}(\mathcal{F}, \mathcal{G})) \stackrel{\sim}{\longrightarrow} \operatorname{Hom}(\mathcal{F} \times \mathcal{H}, \mathcal{G})
\]
(对其中出现的三个自变量函子性的同构)，允许公理化地、以范畴 Cat 的內在术语解释 Hom$(\mathcal{F}, \mathcal{G})$，使得 Hom 范畴的已知公式表看起来是"Hom 对象"存在的范畴 (如 Cat) 中有效公式表的特殊情形. Hom$_{\mathcal{E}/-}(\mathcal{F}, \mathcal{G})$ 有类似的解释，当重新假设 $\mathcal{E}$ 任意时，由公式
\[
\operatorname{Hom}(\mathcal{H}, \mathbf{Hom}_{\mathcal{E}/-}(\mathcal{F}, \mathcal{G})) \stackrel{\sim}{\longrightarrow} \operatorname{Hom}_{\mathcal{E}}(\mathcal{F} \times \mathcal{H}, \mathcal{G})
\]
给出 (对三个自变量函子性的同构). 这样，第 2、3 号中阐述的形式性质是更一般结果的特殊情形，后者在"Hom 对象" $_{\mathcal{E}/-}(\mathcal{F}, \mathcal{G})$ (当 $\mathcal{F}, \mathcal{G}$ 是第三个范畴 $\mathcal{E}$ 之上的两个对象时) 存在的范畴中成立.
\end{remarkstar}

\section{纤维范畴；\texorpdfstring{$\mathcal{E}$}{E}-范畴的等价}

设 $\mathcal{F}$ 是 $\mathcal{E}$ 之上的范畴，$S \in \mathrm{Ob}(\mathcal{E})$. 称 $\mathcal{F}$ 在 $S$ 处的\emph{纤维范畴} (fiber category)\index[term]{xianweifanchou (xianwei)@纤维范畴 (纤维)} 为 $\mathcal{F}$ 的子范畴 $\mathcal{F}_S$，它是 $\mathcal{E}$ 的由 $S$ 定义的单点子范畴的原像. 故 $\mathcal{F}_S$ 的对象是 $p(\xi) = S$ 的 $\mathcal{F}$ 的对象 $\xi$，其态射是 $p(u) = \mathrm{id}_S$ 的 $\mathcal{F}$ 的态射 $u$，i.e. $\mathcal{F}$ 中的 $S$-态射. 当然，$\mathcal{F}_S$ 典范同构于纤维积 $\mathcal{F} \times_{\mathcal{E}} \{S\}$，其中 $\{S\}$ 表示 $\mathcal{E}$ 的由 $S$ 定义的单点子范畴，配备其到 $\mathcal{E}$ 的包含函子. (考虑到基变换的传递性) 由此得出，若作基变换 $\lambda \colon \mathcal{E}' \to \mathcal{E}$，则对每个 $S' \in \mathrm{Ob}(\mathcal{E}')$，投影 $\mathrm{pr}_1 \colon \mathcal{F}' = \mathcal{F} \times_{\mathcal{E}} \mathcal{E}' \to \mathcal{F}$ 诱导同构
\[
\mathcal{F}'_{S'} \longrightarrow \mathcal{F}_S \quad (\text{其中 } S = \lambda(S')).
\]

\begin{proposition}\label{prop:VI.4.1}
设 $f \colon \mathcal{F} \to \mathcal{G}$ 是 $\mathcal{E}$-函子. 若 $f$ 全忠实，则对任何基变换 $\mathcal{E}' \to \mathcal{E}$，相应函子 $f' \colon \mathcal{F}' = \mathcal{F} \times_{\mathcal{E}} \mathcal{E}' \to \mathcal{G}' = \mathcal{G} \times_{\mathcal{E}} \mathcal{E}'$ 全忠实.
\end{proposition}

验证是直接的；更一般地，可以证明全忠实函子的任何投射极限 (这里 $f$ 和 $\mathcal{E}$，$\mathcal{E}'$ 中的恒等函子) 是全忠实函子.

注意与 \ref{prop:VI.4.1} 类似、把"全忠实"换成"范畴等价"的断言是假的，对 $\mathcal{G} = \mathcal{E}$ 已如此. 然而：

\begin{proposition}\label{prop:VI.4.2}
设 $f \colon \mathcal{F} \to \mathcal{G}$ 是 $\mathcal{E}$-函子. 下列条件等价：
\begin{enumerate}[(i)]
\item 存在 $\mathcal{E}$-函子 $g \colon \mathcal{G} \to \mathcal{F}$ 和 $\mathcal{E}$-同构
\[
gf \stackrel{\sim}{\longrightarrow} \mathrm{id}_{\mathcal{F}}, \quad fg \stackrel{\sim}{\longrightarrow} \mathrm{id}_{\mathcal{G}}.
\]
\item 对 $\mathcal{E}$ 之上的每个范畴 $\mathcal{E}'$，函子
\[
f' = f \times_{\mathcal{E}} \mathcal{E}' \colon \mathcal{F}' = \mathcal{F} \times_{\mathcal{E}} \mathcal{E}' \longrightarrow \mathcal{G}' = \mathcal{G} \times_{\mathcal{E}} \mathcal{E}'
\]
是范畴等价.
\item $f$ 是范畴等价，且对每个 $S \in \mathrm{Ob}(\mathcal{E})$，$f$ 诱导的函子 $f_S \colon \mathcal{F}_S \to \mathcal{G}_S$ 是范畴等价.
\item[(iii bis)] $f$ 全忠实，且对每个 $S \in \mathrm{Ob}(\mathcal{E})$ 和每个 $\eta \in \mathrm{Ob}(\mathcal{G}_S)$，存在 $\xi \in \mathrm{Ob}(\mathcal{F}_S)$ 和 $S$-同构 $u \colon f(\xi) \to \eta$.
\end{enumerate}
\end{proposition}

\paragraph*{证明.} 显然 (i) 蕴含 $f$ 是范畴等价 (该概念用同样条件定义，但其中的函子同构不必是 $\mathcal{E}$-态射). 另一方面，由前一节的函子性得出，条件 (i) 在基变换 $\mathcal{E}' \to \mathcal{E}$ 后保持. 由此得 (i)$\Rightarrow$(ii). 显然 (ii)$\Rightarrow$(iii)，因为只需作 $\mathcal{E}' = \mathcal{E}$ 和 $\mathcal{E}' = \{S\}$. (iii)$\Rightarrow$(iii bis) 更平凡，剩下证明 (iii bis)$\Rightarrow$(i). 为此，对每个 $\eta \in \mathrm{Ob}(\mathcal{G})$ 选择 $g(\eta) \in \mathrm{Ob}(\mathcal{F})$ 和同构 $u(\eta) \colon f(g(\eta)) \to \eta$，使得 $q(u(\eta)) = \mathrm{id}_S$，其中 $S = q(\eta)$. 由第二个条件 (iii bis) 这是可能的. 由于 $f$ 全忠实，众所周知且立即看出，$g$ 可以唯一地看作 $\mathcal{G}$ 到 $\mathcal{F}$ 的函子，使诸 $u(\eta)$ 定义函子同态 (从而是同构) $u \colon fg \stackrel{\sim}{\longrightarrow} \mathrm{id}_{\mathcal{G}}$. 此外，由构造 $g$ 是 $\mathcal{E}$-函子，$u$ 是 $\mathcal{E}$-同态. 与前述数据对应的是函子同构 $v \colon gf \to \mathrm{id}_{\mathcal{F}}$，由条件 $f * v = u * f$ 定义，立即看出它也是 $\mathcal{E}$-态射，证毕.

\begin{definition}\label{def:VI.4.3}
若前面的条件满足，则称 $f$ 是 $\mathcal{E}$ 之上的范畴等价，或 $\mathcal{E}$-等价.
\end{definition}

\begin{corollary}\label{cor:VI.4.4}
假设投影函子 $p \colon \mathcal{F} \to \mathcal{E}$ 是可运输函子 (transportable functor)，i.e. 对 $\mathcal{E}$ 中的每个同构 $\alpha \colon T \to S$ 和 $\mathcal{F}_T$ 的每个对象 $\xi$，存在 $\mathcal{F}$ 中以 $\xi$ 为源、满足 $p(u) = \alpha$ 的同构 $u$. 则 $\mathcal{F}$ 到 $\mathcal{G}$ 的每个是范畴等价的 $\mathcal{E}$-函子 $f \colon \mathcal{F} \to \mathcal{G}$ 是 $\mathcal{E}$-等价.
\end{corollary}

由判则 (iii bis) 得出.

\begin{corollary}\label{cor:VI.4.5}
设 $f \colon \mathcal{F} \to \mathcal{G}$ 是 $\mathcal{E}$-等价. 则对 $\mathcal{E}$ 之上的每个范畴 $\mathcal{H}$，相应函子
\[
\operatorname{Hom}_{\mathcal{E}/-}(\mathcal{G}, \mathcal{H}) \longrightarrow \operatorname{Hom}_{\mathcal{E}/-}(\mathcal{F}, \mathcal{H})
\]
\[
\mathbf{Hom}_{\mathcal{E}/-}(\mathcal{H}, \mathcal{F}) \longrightarrow \mathbf{Hom}_{\mathcal{E}/-}(\mathcal{H}, \mathcal{G})
\]
(参见第 2 号) 是范畴等价.
\end{corollary}

由判则 (i) 按通常推理得出.

\section{笛卡儿态射、逆像、笛卡儿函子}

设 $\mathcal{F}$ 是 $\mathcal{E}$ 之上的范畴，投影函子为 $p$.

\begin{definition}\label{def:VI.5.1}
考虑 $\mathcal{F}$ 中的态射
\[
\alpha \colon \eta \longrightarrow \xi
\]
并设
\[
S = p(\xi), \quad T = p(\eta), \quad f = p(\alpha).
\]
称 $\alpha$ 是\emph{笛卡儿态射} (cartesian morphism)\index[term]{dikaertaishe@笛卡儿态射}，若对每个 $\eta' \in \mathrm{Ob}(\mathcal{F}_T)$ 和每个 $f$-态射 $u \colon \eta' \to \xi$，存在唯一 $T$-态射 $\overline{u} \colon \eta' \to \eta$，使得 $u = \alpha \circ \overline{u}$.

因此，这意味着对每个 $\eta' \in \mathrm{Ob}(\mathcal{F}_T)$，映射 $v \mapsto \alpha \circ v$
\begin{equation}
\operatorname{Hom}_T(\eta', \eta) \longrightarrow \operatorname{Hom}_f(\eta', \xi) \tag{i}
\label{eq:VI.i3}
\end{equation}
双射. 这也意味着对 $(\eta, \alpha)$ 表示第二个成员关于 $\eta'$ 的函子 $\mathcal{F}_T^{\circ} \to$ Ens. 若对 $\mathcal{E}$ 中给定的态射 $f \colon T \to S$ 和给定的 $\xi \in \mathrm{Ob}(\mathcal{F}_S)$，存在这样的对 $(\eta, \alpha)$，i.e. 存在 $\mathcal{F}$ 中以 $\xi$ 为靶、满足 $p(\alpha) = f$ 的笛卡儿态射 $\alpha$，则 $\eta$ 在 $\mathcal{F}_T$ 中在唯一同构意义下确定. 此时称 $\xi$ 经 $f$ 的逆像 (inverse image) 存在，$\mathcal{F}_T$ 中配备 $f$-笛卡儿态射 $\alpha \colon \eta \to \xi$ 的对象 $\eta$ 称为 $\xi$ 经 $f$ 的逆像. 常假设在每次存在时选定了这样的逆像 ($\mathcal{F}$ 固定)；逆像用 $f_{\mathcal{F}}^{*}(\xi)$，或简写 $f^{*}(\xi)$，或当这些记号不引起混淆时用 $\xi \times_S T$ 表示；典范态射 $\alpha \colon \eta \to \xi$ 在下面用 $\alpha_f(\xi)$ 表示. 若对每个 $\xi \in \mathrm{Ob}(\mathcal{F}_S)$，$\xi$ 经 $f$ 的逆像存在，也说 $\mathcal{F}$ 中经 $f$ 的逆像函子存在，$f^{*}(\xi)$ 于是成为 $\xi$ 的共变函子，从 $\mathcal{F}_S$ 到 $\mathcal{F}_T$. 这来自 (i) 中第二个成员对 $\xi$ 的共变依赖，i.e. 精确地说表示 $\mathcal{F}_T^{\circ} \times \mathcal{F}_S$ 到 Ens 的函子. 对 $f^{*}(\xi)$ 的这种函子性依赖如下明确：考虑 $f$-笛卡儿态射
\[
\alpha \colon \eta \longrightarrow \xi, \quad \alpha' \colon \eta' \longrightarrow \xi'
\]
和 $S$-态射 $\lambda \colon \xi \to \xi'$，则存在且只存在一个 $T$-态射 $\mu \colon \eta \to \eta'$，使得
\[
\alpha'\mu = \lambda\alpha
\]
(由 $\alpha'$ 笛卡儿得出). 还注意如下立即事实：考虑 $\mathcal{F}$ 中的交换图式
\[
\begin{tikzcd}[row sep=large, column sep=large]
\xi \arrow[d,"\lambda"] & \eta \arrow[l,"\alpha"] \arrow[d,"\mu"] \\
\xi' & \eta' \arrow[l,"\alpha'"]
\end{tikzcd}
\]
其中 $\alpha$ 和 $\alpha'$ 是 $f$-态射，$\lambda$ 是 $S$-同构，$\mu$ 是 $T$-同构. 为使 $\alpha$ 笛卡儿，必要且充分条件是 $\alpha'$ 笛卡儿.
\end{definition}

\begin{definition}\label{def:VI.5.2}
$\mathcal{E}$-函子 $F \colon \mathcal{F} \to \mathcal{G}$ 称为\emph{笛卡儿函子} (cartesian functor)\index[term]{dikaerhanzi@笛卡儿函子}，若它把笛卡儿态射变为笛卡儿态射. 用 $\mathbf{Hom}_{\mathrm{cart}}(\mathcal{F}, \mathcal{G})$\index[notations]{$\mathbf{Hom}_{\mathrm{cart}}(\mathcal{F}, \mathcal{G})$}\index[notations]{$\mathbf{Hom}_{\mathrm{cart}}(\mathcal{F}, \mathcal{G})$} 记 $\mathbf{Hom}_{\mathcal{E}/-}(\mathcal{F}, \mathcal{G})$ 的由笛卡儿函子组成的全子范畴.
\end{definition}

例如，借助恒等函子把 $\mathcal{E}$ 看作 $\mathcal{E}$ 之上的范畴，$\mathcal{E}$ 的每个态射都是笛卡儿态射，因此 $\mathcal{E}$ 到 $\mathcal{F}$ 的笛卡儿函子是把 $\mathcal{E}$ 的每个态射变为笛卡儿态射的截面函子 $F \colon \mathcal{E} \to \mathcal{F}$；这样的函子称为 $\mathcal{F}$ 在 $\mathcal{E}$ 上的笛卡儿截面 (cartesian section).

\begin{proposition}\label{prop:VI.5.3}
(i) 是 $\mathcal{E}$-等价的函子 $F \colon \mathcal{F} \to \mathcal{G}$ 是笛卡儿函子. (ii) 设 $F$，$G$ 是两个同构的 $\mathcal{E}$-函子 $\mathcal{F} \to \mathcal{G}$. 若其中之一笛卡儿，则另一个亦然. (iii) 两个笛卡儿函子 $\mathcal{F} \to \mathcal{G}$ 和 $\mathcal{G} \to \mathcal{H}$ 的复合是笛卡儿函子.
\end{proposition}

断言 (iii) 由定义平凡，(ii) 由 \ref{def:VI.5.2} 前的评注得出，(i) 容易由定义和判则 \ref{prop:VI.4.2} (iii) 得出；更精确地说，$\mathcal{F}$ 中的态射 $\alpha$ 笛卡儿当且仅当 $F(\alpha)$ 笛卡儿.

\begin{corollary}\label{cor:VI.5.4}
设 $F \colon \mathcal{F} \to \mathcal{G}$ 是 $\mathcal{E}$-等价. 则对 $\mathcal{E}$ 之上的每个范畴 $\mathcal{H}$，相应函子 $G \mapsto G \circ F$ 和 $G \mapsto F \circ G$ 诱导范畴等价：
\[
\begin{array}{rl}
& \mathbf{Hom}_{\mathrm{cart}}(\mathcal{G}, \mathcal{H}) \xrightarrow{\approx} \mathbf{Hom}_{\mathrm{cart}}(\mathcal{F}, \mathcal{H}) \\
& \mathbf{Hom}_{\mathrm{cart}}(\mathcal{H}, \mathcal{F}) \xrightarrow{\approx} \mathbf{Hom}_{\mathrm{cart}}(\mathcal{H}, \mathcal{G})
\end{array}
\]
\end{corollary}

由 \ref{prop:VI.4.2} 判则 (i) 和 \ref{prop:VI.5.3} (i) (ii) (iii) 按通常方式推出. 可以精确化：$\mathcal{E}$-函子 $G \colon \mathcal{G} \to \mathcal{H}$ 笛卡儿当且仅当 $G \circ F$ 笛卡儿，同样 $\mathcal{E}$-函子 $G \colon \mathcal{H} \to \mathcal{F}$ 笛卡儿当且仅当 $F \circ G$ 笛卡儿.

由 \ref{cor:VI.5.4} (iii) 得出，若考虑 $\mathbf{Cat}_{/\mathcal{E}}$ 的子范畴 $\mathbf{Cat}_{/\mathcal{E}}^{\mathrm{cart}}$ (其对象与 $\mathbf{Cat}_{/\mathcal{E}}$ 的相同，态射是笛卡儿函子)，则如第 2 号那样有配对：
\[
\operatorname{Hom}_{\mathrm{cart}}(\mathcal{F}, \mathcal{G}) \times \operatorname{Hom}_{\mathrm{cart}}(\mathcal{G}, \mathcal{H}) \longrightarrow \operatorname{Hom}_{\mathrm{cart}}(\mathcal{F}, \mathcal{H})
\]
由第 2 号的配对诱导，使 $\mathbf{Hom}_{\mathrm{cart}}(\mathcal{F}, \mathcal{G})$ 可以看作 $\mathcal{F}, \mathcal{G}$ 的函子，从范畴 $(\mathbf{Cat}_{/\mathcal{E}}^{\mathrm{cart}})^{\circ} \times \mathbf{Cat}_{/\mathcal{E}}^{\mathrm{cart}}$ 到 Cat. 我们主要对 $\mathcal{F} = \mathcal{G}$ 的情形需要这个评注.

\begin{definition}\label{def:VI.5.5}
设 $\mathcal{F}$ 是 $\mathcal{E}$ 之上的范畴. 用
\[
\varprojlim \mathcal{F}/\mathcal{E}\index[notations]{$\varprojlim \mathcal{F}/\mathcal{E}$}
\]
记 $\mathcal{E}$ 到 $\mathcal{F}$ 的笛卡儿 $\mathcal{E}$-函子的范畴，i.e. $\mathcal{F}$ 在 $\mathcal{E}$ 上的笛卡儿截面的范畴.
\end{definition}

由前面所述，Lim $\mathcal{F}/\mathcal{E}$ 是 $\mathcal{F}$ 的函子，从范畴 Cat$^{\mathrm{cart}}_{/\mathcal{E}}$ 到 Cat.

下面将看到这个运算 Lim 与范畴投射极限概念的关系，以及大量例子.

\section{纤维范畴与预纤维范畴。其中的积与基变换}

\begin{definition}\label{def:VI.6.1}
$\mathcal{E}$ 之上的范畴 $\mathcal{F}$ 称为纤维范畴 (fibered category)\index[term]{xianweifanchou@纤维范畴} (此时说函子 $\mathcal{F} \to \mathcal{E}$ 是纤维的 (fibrant))，若它满足下列两个公理：
\begin{itemize}\tightlist
\item[Fib$_1$] 对 $\mathcal{E}$ 中的每个态射 $f \colon T \to S$，$\mathcal{F}$ 中经 $f$ 的逆像函子存在.
\item[Fib$_{II}$] 两个笛卡儿态射的复合是笛卡儿态射.
\end{itemize}
$\mathcal{E}$ 之上满足条件 Fib$_1$ 的范畴 $\mathcal{F}$ 称为 $\mathcal{E}$ 上的\emph{预纤维范畴} (prefibered category)\index[term]{yuxianweifanchou@预纤维范畴}.

若 $\mathcal{F}$ 是 $\mathcal{E}$ 上的纤维范畴 (相应地：预纤维范畴)，$\mathcal{F}$ 的子范畴 $\mathcal{G}$ 称为纤维子范畴 (相应地：预纤维子范畴)，若它是 $\mathcal{E}$ 上的纤维范畴 (相应地：预纤维范畴)，且包含函子笛卡儿. 例如若 $\mathcal{G}$ 是 $\mathcal{F}$ 的全子范畴，看到这意味着：对 $\mathcal{E}$ 中的每个态射 $f \colon T \to S$ 和每个 $\xi \in \mathrm{Ob}(\mathcal{G}_S)$，$f_{\mathcal{F}}^{*}(\xi)$ 与 $\mathcal{G}_T$ 的某个对象 $T$-同构. 另一个有趣的情形是：$\mathcal{F}$ 是 $\mathcal{E}$ 上的纤维范畴，考虑 $\mathcal{F}$ 的子范畴 $\mathcal{G}$，其对象相同，态射是 $\mathcal{F}$ 的笛卡儿态射；特别 $\mathcal{G}_S$ 的态射是 $\mathcal{F}_S$ 的同构. 立即看到这确实是 $\mathcal{F}$ 的纤维子范畴，因为在与 $\mathcal{F}$ 中 $f$-笛卡儿态射 $\alpha$ 相关的双射
\[
\operatorname{Hom}_T(\eta', \eta) \xrightarrow{\sim} \operatorname{Hom}_f(\eta', \xi)
\]
中，与第一成员的 $T$-同构对应的是第二成员的笛卡儿态射. 按定义，$\mathcal{E}$ 到 $\mathcal{F}$ 的笛卡儿截面于是与 $\mathcal{E}$ 到 $\mathcal{G}$ 的任意 $\mathcal{E}$-函子一一对应 (但注意自然函子
\[
\mathbf{Hom}_{\mathcal{E}/-}(\mathcal{E}, \mathcal{G}) \longrightarrow \mathbf{Hom}_{\mathrm{cart}}(\mathcal{E}, \mathcal{F}) = \underline{\mathrm{Lim}}(\mathcal{F}/\mathcal{E})
\]
忠实，但一般不严格全忠实，i.e. 不是同构).
\end{definition}

\begin{remarkstar}
设 $\mathcal{F}$ 是 $\mathcal{E}$ 之上的范畴. 下列条件等价：(i) $\mathcal{F}$ 的所有态射笛卡儿；(ii) $\mathcal{F}$ 是 $\mathcal{E}$ 上的纤维范畴，且诸 $\mathcal{F}_S$ 是群胚 (groupoid)\index[term]{qunpei@群胚} (i.e. $\mathcal{F}_S$ 中每个态射是同构). 此时称 $\mathcal{F}$ 是 $\mathcal{E}$ 上的\emph{群胚纤维范畴} (fibered category in groupoids)\index[term]{qunpeixianweifanchou@群胚纤维范畴}. 在"模理论"中遇到的正是它们. 若 $\mathcal{E}$ 是群胚，证明条件 (i) 和 (ii) 还等价于：(iii) $\mathcal{F}$ 是群胚，且投影函子 $p \colon \mathcal{F} \to \mathcal{E}$ 可运输 (参见 \ref{cor:VI.4.4}). 例如，若 $\mathcal{E}$ 和 $\mathcal{F}$ 是使 $\mathrm{Ob}(\mathcal{E})$ 和 $\mathrm{Ob}(\mathcal{F})$ 约化为一点的群胚，从而 $\mathcal{E}$ 和 $\mathcal{F}$ 在同构意义下由群 $E$ 和 $F$ 定义，且函子 $p \colon \mathcal{F} \to \mathcal{E}$ 由群同态 $p \colon F \to E$ 定义，则 $\mathcal{F}$ 在 $\mathcal{E}$ 上纤维当且仅当 $p$ 满射，i.e. 当且仅当 $p$ 定义群 $E$ 以群 $G = \operatorname{Ker} p$ 为核的扩张.
\end{remarkstar}

\begin{proposition}\label{prop:VI.6.2}
设 $F \colon \mathcal{F} \to \mathcal{G}$ 是 $\mathcal{E}$-等价. 为使 $\mathcal{F}$ 是 $\mathcal{E}$ 上的纤维范畴 (相应地：预纤维范畴)，必要且充分条件是 $\mathcal{G}$ 如此.
\end{proposition}

由定义和上面指出的评注 (即 $\mathcal{F}$ 中的态射 $\alpha$ 笛卡儿当且仅当 $F(\alpha)$ 笛卡儿) 容易得出.

\begin{proposition}\label{prop:VI.6.3}
设 $\mathcal{F}_1$，$\mathcal{F}_2$ 是 $\mathcal{E}$ 之上的两个范畴，$\alpha = (\alpha_1, \alpha_2)$ 是 $\mathcal{F} = \mathcal{F}_1 \times_{\mathcal{E}} \mathcal{F}_2$ 中的态射. 为使 $\alpha$ 笛卡儿，必要且充分条件是它的两个分量笛卡儿.
\end{proposition}

确实设 $\xi_i$ 是 $\alpha_i$ 的靶，$\eta_i$ 是源，$f \colon T \to S$ 是使 $\alpha_1$ 和 $\alpha_2$ 成为 $f$-态射的 $\mathcal{E}$ 的态射. 对每个 $\eta' = (\eta'_1, \eta'_2) \in \mathcal{F}_T$，有交换图式 (其中垂直箭是双射). 故若两个水平箭之一是双射，则另一个亦然. 这已表明：若 $\alpha_1$，$\alpha_2$ 笛卡儿 (从而第二个水平箭双射)，则 $\alpha$ 笛卡儿. 反过来，在上图式中取 $\eta'_i = \eta_i$，由此得 Hom$_T(\eta'_i, \eta_i) \neq \emptyset$，先对 $i = 2$，这证明 $\alpha_1$ 笛卡儿，再对 $i = 1$，证明 $\alpha_2$ 笛卡儿.

\begin{corollary}\label{cor:VI.6.4}
设 $\mathcal{F} = \mathcal{F}_1 \times_{\mathcal{E}} \mathcal{F}_2$，$F = (F_1, F_2)$ 是 $\mathcal{E}$-函子 $\mathcal{G} \to \mathcal{F}$. 为使 $F$ 笛卡儿，必要且充分条件是 $F_1$ 和 $F_2$ 笛卡儿. 由此得到范畴同构
\[
\mathbf{Hom}_{\mathrm{cart}}(\mathcal{G}, \mathcal{F}_1 \times_{\mathcal{E}} \mathcal{F}_2) \xrightarrow{\sim} \mathbf{Hom}_{\mathrm{cart}}(\mathcal{G}, \mathcal{F}_1) \times \mathbf{Hom}_{\mathrm{cart}}(\mathcal{G}, \mathcal{F}_2)
\]
特别 (取 $\mathcal{G} = \mathcal{E}$) 得到范畴同构
\[
\varprojlim (\mathcal{F}_1 \times_{\mathcal{E}} \mathcal{F}_2/\mathcal{E}) \xrightarrow{\sim} \varprojlim (\mathcal{F}_1/\mathcal{E}) \times \varprojlim (\mathcal{F}_2/\mathcal{E})
\]
\end{corollary}

\begin{corollary}\label{cor:VI.6.5}
设 $\mathcal{F}_1$ 和 $\mathcal{F}_2$ 是 $\mathcal{E}$ 上的两个纤维范畴 (相应地：预纤维范畴)，则它们的纤维积 $\mathcal{F} = \mathcal{F}_1 \times_{\mathcal{E}} \mathcal{F}_2$ 是 $\mathcal{E}$ 上的纤维范畴 (相应地：预纤维范畴).
\end{corollary}

这些结果还推广到 $\mathcal{E}$ 之上任意族范畴的纤维积的情形.

\begin{proposition}\label{prop:VI.6.6}
设 $\mathcal{F}$ 是 $\mathcal{E}$ 之上的范畴，投影函子为 $p$，$\lambda \colon \mathcal{E}' \to \mathcal{E}$ 是函子，借助投影函子 $p' = p \times_{\mathcal{E}} \mathrm{id}_{\mathcal{E}'}$ 把 $\mathcal{F}' = \mathcal{F} \times_{\mathcal{E}} \mathcal{E}'$ 看作 $\mathcal{E}'$ 之上的范畴. 设 $\alpha'$ 是 $\mathcal{F}'$ 的态射，为使 $\alpha'$ 是笛卡儿态射，必要且充分条件是它在 $\mathcal{F}$ 中的像 $\alpha$ 笛卡儿.
\end{proposition}

证明是直接的，留给读者.

\begin{corollary}\label{cor:VI.6.7}
对 $\mathcal{E}$ 之上范畴的每个笛卡儿函子 $F \colon \mathcal{F} \to \mathcal{G}$，函子 $F' = F \times_{\mathcal{E}} \mathcal{E}'$ (从 $\mathcal{F}' = \mathcal{F} \times_{\mathcal{E}} \mathcal{E}'$ 到 $\mathcal{G}' = \mathcal{G} \times_{\mathcal{E}} \mathcal{E}'$) 笛卡儿.
\end{corollary}

因此，第 3 号所考虑的函子 $\mathbf{Hom}_{\mathcal{E}}(\mathcal{F}, \mathcal{G}) \to \mathbf{Hom}_{\mathcal{E}'}(\mathcal{F}', \mathcal{G}')$ 诱导函子
\[
\operatorname{Hom}_{\mathrm{cart}}(\mathcal{F}, \mathcal{G}) \longrightarrow \operatorname{Hom}_{\mathrm{cart}}(\mathcal{F}', \mathcal{G}');
\]
换言之，对固定的 $\mathcal{F}, \mathcal{G}$，可以把
\[
\mathbf{Hom}_{\mathrm{cart}}(\mathcal{F} \times_{\mathcal{E}} \mathcal{E}', \mathcal{G} \times_{\mathcal{E}} \mathcal{E}')
\]
看作 $\mathcal{E}'$ 的函子，从范畴 $\mathbf{Cat}_{/\mathcal{E}}^{\circ}$ 到 Cat. 若也让 $\mathcal{F}, \mathcal{G}$ 变动，得到从范畴 $\mathbf{Cat}_{/\mathcal{E}}^{\circ} \times (\mathbf{Cat}_{/\mathcal{E}}^{\mathrm{cart}})^{\circ} \times \mathbf{Cat}_{/\mathcal{E}}^{\mathrm{cart}}$ 到 Cat 的函子. 考虑到第 3 号所考虑的同构
\[
\operatorname{Hom}_{\mathcal{E}'}(\mathcal{F}', \mathcal{G}') \stackrel{\sim}{\longrightarrow} \operatorname{Hom}_{\mathcal{E}}(\mathcal{F} \times_{\mathcal{E}} \mathcal{E}', \mathcal{G}),
\]
则 $\mathcal{F}'$ 到 $\mathcal{G}'$ 的笛卡儿 $\mathcal{E}'$-函子与把"第一投影是 $\mathcal{F}$ 的笛卡儿态射"的每个态射变为 $\mathcal{G}$ 的笛卡儿态射的 $\mathcal{E}$-函子 $\mathcal{F} \times_{\mathcal{E}} \mathcal{E}' \to \mathcal{G}$ 对应. 取 $\mathcal{F} = \mathcal{E}$，(更换记号后) 得到：

\begin{corollary}\label{cor:VI.6.8}
Lim$(\mathcal{F}'/\mathcal{E}')$ 同构于 $\mathbf{Hom}_{\mathcal{E}/-}(\mathcal{E}', \mathcal{F})$ 的由把任意态射变为笛卡儿态射的 $\mathcal{E}$-函子 $\mathcal{E}' \to \mathcal{F}$ 组成的全子范畴. 特别，若 $\mathcal{F}$ 是纤维范畴且 $\widetilde{\mathcal{F}}$ 是以 $\mathcal{F}$ 的笛卡儿态射为态射的 $\mathcal{F}$ 的子范畴，则双射
\[
\mathrm{Ob} \varprojlim (\mathcal{F}'/\mathcal{E}') \xrightarrow{\sim} \mathrm{Hom}_{\mathcal{E}/-}(\mathcal{E}', \widetilde{\mathcal{F}}).
\]
这精确了 $\underline{\mathrm{Lim}}(\mathcal{F} \times_{\mathcal{E}} \mathcal{E}'/\mathcal{E}')$ 应被看作 $\mathcal{E}'$ 和 $\mathcal{F}$ 的函子的方式，从范畴 Cat$^{\circ}_{/\mathcal{E}} \times \mathbf{Cat}^{\mathrm{cart}}_{/\mathcal{E}}$ 到 Cat. 后面在 $\mathcal{F}$ 限制为纤维范畴时将看到对 $\mathcal{E}'$ 的更完全的函子依赖.
\end{corollary}

\begin{corollary}\label{cor:VI.6.9}
设 $\mathcal{F}$ 是 $\mathcal{E}$ 上的纤维范畴 (相应地：预纤维范畴)，则 $\mathcal{F}' = \mathcal{F} \times_{\mathcal{E}} \mathcal{E}'$ 是 $\mathcal{E}'$ 上的纤维范畴 (相应地：预纤维范畴).
\end{corollary}

\begin{proposition}\label{prop:VI.6.10}
设 $\mathcal{F}, \mathcal{G}$ 是 $\mathcal{E}$ 上的预纤维范畴，$F$ 是 $\mathcal{F}$ 到 $\mathcal{G}$ 的笛卡儿 $\mathcal{E}$-函子. 为使 $F$ 忠实 (相应地：全忠实，相应地：是 $\mathcal{E}$-等价)，必要且充分条件是：对每个 $S \in \mathrm{Ob}(\mathcal{E})$，诱导函子 $F_S \colon \mathcal{F}_S \to \mathcal{G}_S$ 忠实 (相应地：全忠实，相应地：是等价).
\end{proposition}

由定义立即证明.

为结束本节，给出纤维范畴的若干性质，使用公理 Fib$_{II}$.

\begin{proposition}\label{prop:VI.6.11}
设 $\mathcal{F}$ 是 $\mathcal{E}$ 上的预纤维范畴. 为使 $\mathcal{F}$ 纤维，必要且充分条件是它满足下列条件：
\begin{itemize}\tightlist
\item[Fib$'_{II}$：] 设 $\alpha \colon \eta \to \xi$ 是 $\mathcal{F}$ 中 $\mathcal{E}$ 的态射 $f \colon T \to S$ 之上的笛卡儿态射. 对 $\mathcal{E}$ 中的每个态射 $g \colon U \to T$ 和每个 $\zeta \in \mathrm{Ob}(\mathcal{F}_U)$，映射 $u \mapsto \alpha \circ u$
\[
\mathrm{Hom}_g(\zeta, \eta) \longrightarrow \mathrm{Hom}_{fg}(\zeta, \xi)
\]
双射.
\end{itemize}
\end{proposition}

换言之，在 $\mathcal{E}$ 上的纤维范畴中，笛卡儿图式由先验更强的性质刻画 (在下列陈述中取 $g = \mathrm{id}_T$ 即得定义中的性质).

\begin{corollary}\label{cor:VI.6.12}
设 $\mathcal{F}$ 是 $\mathcal{E}$ 之上的范畴，$\alpha$ 是 $\mathcal{F}$ 中的态射. 为使 $\alpha$ 是同构，必须 $p(\alpha) = f$ 是同构且 $\alpha$ 笛卡儿；若 $\mathcal{F}$ 在 $\mathcal{E}$ 上纤维，则逆命题成立.
\end{corollary}

事实上，若 $\alpha$ 是同构则 $f = p(\alpha)$ 显然亦然；对每个 $\eta' \in \mathrm{Ob}(\mathcal{F}_T)$，映射 $u \mapsto \alpha \circ u$
\[
\operatorname{Hom}(\eta', \eta) \longrightarrow \operatorname{Hom}(\eta', \xi)
\]
双射；由于 $f$ 是同构，立即看到第一成员的 $T$-态射当且仅当其像为第二成员的 $f$-态射，故得到双射
\[
\operatorname{Hom}_T(\eta', \eta) \longrightarrow \operatorname{Hom}_f(\eta', \xi)
\]
证明第一个断言. 反过来，假设 $f$ 是同构且 $\alpha$ 满足 Fib$'_{II}$ 中陈述的条件 (当 $\mathcal{F}$ 在 $\mathcal{E}$ 上纤维时，这意味着 $\alpha$ 笛卡儿)，立即看到对每个 $\zeta \in \mathrm{Ob} \mathcal{F}$，映射 $u \mapsto \alpha \circ u$ (从 Hom$(\zeta, \eta)$ 到 Hom$(\zeta, \xi)$) 双射，故 $\alpha$ 是同构.

\begin{corollary}\label{cor:VI.6.13}
设 $\alpha \colon \eta \to \xi$ 和 $\beta \colon \zeta \to \eta$ 是纤维范畴 $\mathcal{F}$ 中的两个可复合态射. 若 $\alpha$ 笛卡儿，则 $\beta$ 笛卡儿当且仅当 $\alpha\beta$ 笛卡儿.
\end{corollary}

用 \ref{prop:VI.6.11} 的强化形式的笛卡儿态射定义.

\section{\texorpdfstring{$\mathcal{E}$}{E} 上的劈开范畴}

\begin{definition}\label{def:VI.7.1}
设 $\mathcal{F}$ 是 $\mathcal{E}$ 之上的范畴. 称 $\mathcal{F}$ 在 $\mathcal{E}$ 上的\emph{劈开} (cleavage)\index[term]{pikai@劈开} 为给每个 $f \in \mathrm{Fl}(\mathcal{E})$ 配 $\mathcal{F}$ 中经 $f$ 的逆像函子 $f^{*}$\index[notations]{$f^{*}$} 的函数. 若 $f = \mathrm{id}_S$ 蕴含 $f^{*} = \mathrm{id}_{\mathcal{F}_S}$，则称该劈开规范化 (normalized)\index[term]{guifanhuapikai@规范化劈开}. 配备劈开 (相应地：\emph{规范化劈开}) 的 $\mathcal{E}$ 之上的范畴 $\mathcal{F}$ 称为\emph{劈开范畴} (cleaved category)\index[term]{pikaifanchou@劈开范畴} (相应地：规范化劈开范畴).
\end{definition}

显然，$\mathcal{F}$ 有劈开当且仅当 $\mathcal{F}$ 是 $\mathcal{E}$ 上的预纤维范畴，且此时 $\mathcal{F}$ 有规范化劈开. $\mathcal{F}$ 上的劈开集与 $\mathrm{Fl}(\mathcal{F})$ 的满足下列条件的子集 $K$ 一一对应：
\begin{enumerate}[a)]
\item $K$ 的诸 $\alpha$ 是笛卡儿态射.
\item 对 $\mathcal{E}$ 中的每个态射 $f \colon T \to S$ 和每个 $\xi \in \mathrm{Ob}(\mathcal{F}_S)$，存在 $K$ 中以 $\xi$ 为靶的唯一 $f$-态射.
\end{enumerate}
为使 $K$ 定义的劈开规范化，必要且充分条件是 $K$ 还满足条件
\begin{enumerate}[a)]
\setcounter{enumi}{2}
\item $\mathcal{F}$ 中的恒等态射属于 $K$.
\end{enumerate}
$K$ 的态射可称为所考虑劈开的"\emph{运输态射}" (transport morphisms)\index[term]{yunshutaishe@运输态射}.

$\mathcal{E}$ 上劈开范畴的同构概念是清楚的. 更一般地，可以把 $\mathcal{E}$-劈开范畴的态射定义为把运输态射变为运输态射的 $\mathcal{E}$-范畴函子 $\mathcal{F} \to \mathcal{G}$. (它们特别是笛卡儿函子). 这样，$\mathcal{E}$ 上的劈开范畴是某个范畴的对象，即 $\mathcal{E}$ 上的劈开范畴范畴. 读者将明确积的存在性，它与下列事实有关：若 $\mathcal{E}$ 之上的范畴 $\mathcal{F}$ 是配备各自劈开的 $\mathcal{E}$ 之上范畴 $\mathcal{F}_i$ 的积，则 $\mathcal{F}$ 配备相应的自然劈开. 也把劈开范畴中的基变换概念留给读者明确.

用 $\alpha_f(\xi)$ 记典范态射\index[notations]{$\alpha_f(\xi)$}
\[
\alpha_f(\xi) \colon f^{*}(\xi) \longrightarrow \xi.
\]
已说过，它对 $\xi$ 函子性，i.e. 有函子同态
\[
\alpha_f \colon i_T f^{*} \longrightarrow i_S,
\]
其中对每个 $S \in \mathrm{Ob}(\mathcal{E})$，$i_S$ 记包含函子
\[
i_S \colon \mathcal{F}_S \longrightarrow \mathcal{F}.
\]
现在考虑 $\mathcal{E}$ 中的态射
\[
f \colon T \longrightarrow S \quad \text{和} \quad g \colon U \longrightarrow T,
\]
设 $\xi \in \mathrm{Ob}(\mathcal{F}_S)$，则存在唯一 $U$-态射
\[
c_{f,g}(\xi) \colon g^{*} f^{*}(\xi) \longrightarrow (fg)^{*}(\xi)
\]
使图式交换 (由 $(fg)^{*}(\xi)$ 的定义)：
\begin{equation}
\begin{tikzcd}[row sep=large, column sep=large]
f^{*}(\xi) \arrow[d,"{\alpha_f(\xi)}"] & g^{*}(f^{*}(\xi)) \arrow[l,"{\alpha_g(f^{*}(\xi))}"] \arrow[d,"{c_{f,g}(\xi)}"] \\
\xi & (fg)^{*}(\xi) \arrow[l,"{\alpha_{fg}(\xi)}"]
\end{tikzcd} \tag{172}
\label{eq:VI.172}
\end{equation}
对变动的 $\xi$，这个同态函子性，i.e. 有函子 $\mathcal{F}_S \to \mathcal{F}_U$ 的同态
\[
c_{f,g} \colon g^{*} f^{*} \longrightarrow (fg)^{*}
\]
立即注意：

\begin{proposition}\label{prop:VI.7.2}
为使 $\mathcal{E}$ 上的劈开范畴 $\mathcal{F}$ 纤维，必要且充分条件是诸 $c_{f,g}$ 是同构.
\end{proposition}

取 $f$ 为同构、$g$ 为其逆，并考虑同构 $c_{f,g}$ 和 $c_{g,f}$，由此得出：

\begin{corollary}\label{cor:VI.7.3}
若 $\mathcal{F}$ 是 $\mathcal{E}$ 上的纤维劈开范畴，则对 $\mathcal{E}$ 中的每个同构 $f \colon T \to S$，$f^{*}$ 是范畴等价 $\mathcal{F}_S \to \mathcal{F}_T$.
\end{corollary}

\begin{proposition}\label{prop:VI.7.4}
设 $\mathcal{F}$ 是 $\mathcal{E}$ 上的劈开范畴. 有
\[
\begin{array}{ll}
A) & \left\{ \begin{array}{l}
c_{f,\mathrm{id}_T}(\xi) = \alpha_{\mathrm{id}_T}(f^{*}(\xi)) \\
c_{\mathrm{id}_S,f}(\xi) = f^{*}(\alpha_{\mathrm{id}_S}(\xi))
\end{array} \right. \\
B) & c_{f,gh}(\xi) \cdot c_{g,h}(f^{*}(\xi)) = c_{fg,h}(\xi) \cdot h^{*}(c_{f,g}(\xi))
\end{array}
\]
(在这些公式中，$f, g, h$ 表示态射
\begin{equation}
V \longrightarrow U \longrightarrow T \longrightarrow S \tag{173}
\label{eq:VI.173}
\end{equation}
$\xi$ 是 $\mathcal{F}_S$ 的对象.)
\end{proposition}

在规范化劈开的情形，第一和第二个关系取更简单的形式
\[
A') \qquad c_{f,\mathrm{id}_T} = \mathrm{id}_{f^{*}}, \quad c_{\mathrm{id}_S,f} = \mathrm{id}_{f^{*}}.
\]
至于第三个，它由图式的交换性可视化
\begin{equation}
\begin{tikzcd}[row sep=large, column sep=large]
h^{*} g^{*} f^{*}(\xi) \arrow[r,"{c_{g,h}(f^{*}(\xi))}"] \arrow[d,"{h^{*}(c_{f,g}(\xi))}"] & (gh)^{*}(f^{*}(\xi)) \arrow[d,"{c_{f,gh}(\xi)}"] \\
h^{*}(fg)^{*}(\xi) \arrow[r,"{c_{fg,h}(\xi)}"] & (fgh)^{*}(\xi)
\end{tikzcd} \tag{D}
\label{eq:VI.D}
\end{equation}
在纤维范畴的情形 (诸 $c_{f,g}$ 是同构)，这个交换性可以直观地表达为：依次使用形如 $c_{f,g}$ 的同构不导致"矛盾的等同". 还可以不写自变量 $\xi$ 地写这个公式，用函子同态的卷积积：
\[
c_{fg,h} \circ (h^{*} * c_{f,g}) = c_{f,gh} \circ (c_{g,h} * f^{*}).
\]
\ref{prop:VI.7.4} 前两个公式的证明平凡，略述第三个的证明. 为此，除方形 (D) 外，考虑同态方形
\begin{equation}
\begin{tikzcd}[row sep=large, column sep=large]
g^{*} f^{*}(\xi) \arrow[r,"{\alpha_g(f^{*}(\xi))}"] \arrow[d,"{c_{f,g}(\xi)}"] & f^{*}(\xi) \arrow[d,"{\alpha_f(\xi)}"] \\
(fg)^{*}(\xi) \arrow[r,"{\alpha_{fg}(\xi)}"] & \xi
\end{tikzcd} \tag{(D')}
\end{equation}
它按 $c_{f,g}(\xi)$ 的定义交换. 考虑把 (D) 的顶点与 (D') 的相应顶点用形如 $\alpha$ 的同态连接得到的图式：
\[
\begin{array}{l}
\alpha_h(g^{*} f^{*}(\xi)), \alpha_{gh}(f^{*}(\xi)) \\
\alpha_h((fg)^{*}(\xi)), \alpha_{fgh}(\xi).
\end{array}
\]
如此得到的立方体的四个侧面也交换：对左侧面，来自 (D) 的左列由 (D') 的左列经 $h$ 作用得到，且 $\alpha_h$ 是函子同态；对其余三个，不过是 (D) 其余三边的运算 $c$ 的定义. 故立方体中除顶面外的五个面交换. 由此得出由 (D) 定义的两个 $(fgh)$-态射 $h^{*} g^{*} f^{*}(\xi) \to (fgh)^{*}(\xi)$ 与 $\alpha_{fgh}(\xi) \colon (fgh)^{*}(\xi) \to \xi$ 的复合相同，故由 $(fgh)^{*}$ 的定义它们相等.

下面限于 $\mathcal{E}$ 上的规范化劈开范畴. 这样的范畴产生下列对象：
\begin{enumerate}[a)]
\item 从 $\mathrm{Ob}(\mathcal{E})$ 到 Cat 的映射 $S \mapsto \mathcal{F}_S$.
\item 给每个 $f \in \mathrm{Fl}(\mathcal{E})$ (源为 $T$，靶为 $S$) 配函子 $f^{*} \colon \mathcal{F}_S \to \mathcal{F}_T$ 的映射 $f \mapsto f^{*}$.
\item 给 $\mathcal{E}$ 的每对箭 $(f, g)$ 配函子同态 $c_{f,g} \colon g^{*} f^{*} \to (fg)^{*}$ 的映射 $(f, g) \mapsto c_{f,g}$.
\end{enumerate}
此外，这些数据满足上面给出的公式 $A')$ 和 B) 所表达的条件. (N.B. 若不限于规范化劈开的情形，就必须引入额外对象，即给 $\mathcal{E}$ 的每个对象 $S$ 配函子同态 $\alpha_S \colon (\mathrm{id}_S)^{*} \to \mathrm{id}_{\mathcal{F}_S}$ 的函数 $S \mapsto \alpha_S$；条件 $A')$ 将被条件 A) 代替.)

下面说明如何用前述对象 (在唯一同构意义下) 重建 $\mathcal{E}$ 上的规范化劈开范畴 $\mathcal{F}$.

\section{由伪函子 \texorpdfstring{$\mathcal{E}^{\circ} \to \mathbf{Cat}$}{E^o→Cat} 定义的劈开范畴}

为简短，称从 $\mathcal{E}^{\circ}$ 到 Cat 的\emph{伪函子} (pseudo-functor)\index[term]{weihanzi@伪函子} (应当说：规范化伪函子) 为像上面 a)，b)，c) 那样、满足条件 $A')$ 和 B) 的数据集. 在前一节，我们把 $\mathcal{E}$ 上的规范化劈开范畴与伪函子 $\mathcal{E}^{\circ} \to \mathbf{Cat}$ 对应；这里指出相反的构造. 大多数细节、以及这些构造确实"互逆"的验证留给读者. 精确地说，应当把伪函子 $\mathcal{E}^{\circ} \to$ Cat 看作某个新范畴的对象，并证明我们的构造在它和前一节定义的 $\mathcal{E}$ 之上劈开范畴范畴之间提供互为拟逆的等价.

令
\[
\mathcal{F}_{\circ} = \coprod_{S \in \operatorname{Ob}(\mathcal{E})} \operatorname{Ob}(\mathcal{F}(S)),
\]
集合 $\mathrm{Ob}(\mathcal{F}(S))$ 的和集 (N.B. 这里用 $\mathcal{F}(S)$ 而不用 $\mathcal{F}_S$ 记\index[notations]{$\mathcal{F}_S$}给定伪函子在 $\mathcal{E}$ 的对象 $S$ 处的值，以避免后面的记号混淆). 故有显然映射：
\[
p_{\circ} \colon \mathcal{F}_{\circ} \longrightarrow \mathrm{Ob} \mathcal{E}.
\]
设
\[
\overline{\xi} = (S, \xi), \quad \overline{\eta} = (T, \eta) \qquad (\text{其中 } \xi \in \operatorname{Ob} \mathcal{F}(S), \eta \in \operatorname{Ob} \mathcal{F}(T))
\]
是 $\mathcal{F}_{\circ}$ 的两个元素，$f \in \operatorname{Hom}(T, S)$，令
\[
h_f(\overline{\eta}, \overline{\xi}) = \operatorname{Hom}_{\mathcal{F}(T)}(\eta, f^{*}(\xi)).
\]
若还有 $\mathcal{E}$ 中的态射 $g \colon U \to T$ 和 $\zeta \in \operatorname{Ob} \mathcal{F}(U)$，定义映射，记为 $(u, v) \mapsto u \circ v$
\[
h_f(\overline{\eta}, \overline{\xi}) \times h_g(\overline{\zeta}, \overline{\eta}) \longrightarrow h_{fg}(\overline{\zeta}, \overline{\xi}),
\]
i.e. 映射
\[
\operatorname{Hom}_{\mathcal{F}(T)}(\eta, f^{*}(\xi)) \times \operatorname{Hom}_{\mathcal{F}(U)}(\zeta, g^{*}(\eta)) \longrightarrow \operatorname{Hom}_{\mathcal{F}(U)}(\zeta, (fg)^{*}(\xi)),
\]
由公式
\[
u \circ v = c_{f,g}(\xi) \cdot g^{*}(u) \cdot v,
\]
i.e. $u \circ v$ 是序列
\[
\zeta \xrightarrow{v} g^{*}(\eta) \xrightarrow{g^{*}(u)} g^{*} f^{*}(\xi) \xrightarrow{c_{f,g}(\xi)} (fg)^{*}(\xi)
\]
的复合. 另一方面令
\[
h(\overline{\eta}, \overline{\xi}) = \coprod_{f \in \operatorname{Hom}(T, S)} h_f(\overline{\eta}, \overline{\xi}),
\]
前述配对定义配对
\[
h(\overline{\eta}, \overline{\xi}) \times h(\overline{\zeta}, \overline{\eta}) \longrightarrow h(\overline{\zeta}, \overline{\xi}),
\]
而 $h(\overline{\eta}, \overline{\xi})$ 的定义蕴含显然映射：
\[
p_{\overline{\eta},\overline{\xi}} \colon h(\overline{\eta}, \overline{\xi}) \longrightarrow \operatorname{Hom}(T, S).
\]
这样说，验证下列各点：
\begin{enumerate}[1)]
\item $h(\overline{\eta}, \overline{\xi})$ 的元素之间的复合结合.
\item 对 $\mathcal{F}_{\circ}$ 中的每个 $\overline{\xi} = (\xi, S)$，考虑 $h_{\mathrm{id}_S}(\overline{\xi}, \overline{\xi}) = \operatorname{Hom}_{\mathcal{F}(S)}(\mathrm{id}_S^{*}(\xi), \xi) = \operatorname{Hom}_{\mathcal{F}(S)}(\xi, \xi)$ 的恒等元素，及其在 $h(\overline{\xi}, \overline{\xi})$ 中的像. 该对象是 $h(\overline{\eta}, \overline{\xi})$ 元素之间复合的左幺元和右幺元.
\end{enumerate}
这已表明得到范畴 $\mathcal{F}$，令
\[
\mathrm{Ob} \mathcal{F} = \mathcal{F}_{\circ}, \quad \mathrm{Fl} \mathcal{F} = \coprod_{\overline{\xi}, \overline{\eta} \in \mathcal{F}_{\circ}} h(\overline{\eta}, \overline{\xi}).
\]
(N.B. 不能简单取 $\mathrm{Fl} \mathcal{F}$ 为诸集合 $h(\overline{\eta}, \overline{\xi})$ 的并，因为它们不必不交). 此外：
\begin{enumerate}[1)]
\setcounter{enumi}{2}
\item 映射 $p_{\circ} \colon \mathrm{Ob} \mathcal{F} \to \mathrm{Ob} \mathcal{E}$ 和 $p_1 = (p_{\overline{\eta},\overline{\xi}}) \colon \mathrm{Fl} \mathcal{F} \to \mathrm{Fl} \mathcal{E}$ 定义函子 $p \colon \mathcal{F} \to \mathcal{E}$. 这样，$\mathcal{F}$ 成为 $\mathcal{E}$ 之上的范畴，且显然映射 $h_f(\overline{\eta}, \overline{\xi}) \to \mathrm{Hom}(\overline{\eta}, \overline{\xi})$ 诱导双射
\end{enumerate}
\[
h_f(\overline{\eta}, \overline{\xi}) \stackrel{\sim}{\longrightarrow} \operatorname{Hom}_f(\overline{\eta}, \overline{\xi}).
\]
\begin{enumerate}[1)]
\setcounter{enumi}{3}
\item 显然映射
\end{enumerate}
\[
\operatorname{Ob} \mathcal{F}(S) \longrightarrow \mathcal{F}_{\circ} = \operatorname{Ob} \mathcal{F}, \quad \operatorname{Fl} \mathcal{F}(S) \longrightarrow \operatorname{Fl} \mathcal{F},
\]
(其中第二个由显然映射
\[
\mathrm{Hom}_{\mathcal{F}(S)}(\xi, \xi') = h_{\mathrm{id}_S}(\overline{\xi}, \overline{\xi}') \longrightarrow \mathrm{Hom}(\overline{\xi}, \overline{\xi}')
\]
定义) 定义同构
\[
i_S \colon \mathcal{F}(S) \stackrel{\sim}{\longrightarrow} \mathcal{F}_S.
\]
\begin{enumerate}[1)]
\setcounter{enumi}{4}
\item 对 $\mathcal{F}$ 的每个对象 $\overline{\xi} = (S, \xi)$ 和 $\mathcal{E}$ 的每个态射 $f \colon T \to S$，考虑 $\mathcal{F}_T$ 的元素 $\overline{\eta} = (T, \eta)$ (其中 $\eta = f^{*}(\xi)$) 和 $\mathrm{Hom}(\overline{\eta}, \overline{\xi})$ 的元素 $\alpha_f(\xi)$，它是 $\mathrm{id}_{f^{*}(\xi)}$ 在映射 Hom$_{\mathcal{F}(T)}(f^{*}(\xi), f^{*}(\xi)) = h_f(\overline{\eta}, \overline{\xi}) \to \mathrm{Hom}_f(\overline{\eta}, \overline{\xi})$ 下的像. 该元素笛卡儿，且 $f = \mathrm{id}_S$ 时是 $\overline{\xi}$ 的恒等态射；换言之，诸 $\alpha_f(\xi)$ 定义 $\mathcal{F}$ 在 $\mathcal{E}$ 上的规范化劈开. 此外，由构造，在函子图式中交换
\end{enumerate}
\[
\begin{tikzcd}[row sep=large, column sep=large]
\mathcal{F}(T) \arrow[r,"{f^{*}_{\mathcal{F}}}"] \arrow[d,"{i_T}"] & \mathcal{F}(S) \arrow[d,"{i_S}"] \\
\mathcal{F}_T \arrow[r,"{f^{*}_{\mathcal{F}}}"] & \mathcal{F}_S
\end{tikzcd}
\]
其中 $f^{*}_{\mathcal{F}}$ 是所考虑的劈开在 $\mathcal{F}$ 上相应的经 $f$ 的逆像函子. 最后：6) 与伪函子一起给出的同态 $c_{f,g}$ 经同构 $i_S$ 变为与 $\mathcal{F}$ 的劈开相关的函子同态 $c_{f,g}$.

只给出 1) 的验证 (如果可能，它比其他验证较不平凡). 只需证明形如 $h_f(\overline{\eta}, \overline{\xi})$ 的集合的元素之间复合的结合性. 因此考虑 $\mathcal{E}$ 中的态射
\[
S \xleftarrow{f} T \xleftarrow{g} U \xleftarrow{h} V
\]
和对象
\[
\xi, \eta, \zeta, \tau
\]
分别在 $\mathcal{F}(S), \mathcal{F}(T), \mathcal{F}(U), \mathcal{F}(V)$ 中，最后考虑元素
\[
\begin{array}{rl}
& u \in h_f(\overline{\eta}, \overline{\xi}) = \mathrm{Hom}_{\mathcal{F}(T)}(\eta, f^{*}(\xi)) \\
& v \in h_g(\overline{\zeta}, \overline{\eta}) = \mathrm{Hom}_{\mathcal{F}(U)}(\zeta, g^{*}(\eta)) \\
& w \in h_h(\overline{\tau}, \overline{\zeta}) = \mathrm{Hom}_{\mathcal{F}(V)}(\tau, h^{*}(\zeta)).
\end{array}
\]
要证明公式
\[
(u \circ v) \circ w = u \circ (v \circ w),
\]
它是 $\mathrm{Hom}_{\mathcal{F}(V)}(\tau, (fgh)^{*}(\xi))$ 中的等式. 按定义，该等式的两个成员按下面图式的上轮廓和下轮廓复合得到：
\[
\tau \xrightarrow{w} h^{*}(\zeta) \xrightarrow{h^{*}(v)} h^{*} g^{*}(\eta) \xrightarrow{h^{*} g^{*}(u)} h^{*} g^{*} f^{*}(\xi) \xrightarrow{h^{*}(c_{f,g}(\xi))} h^{*}(fg)^{*}(\xi) \xrightarrow{c_{f,gh}(\xi)} (fgh)^{*}(\xi)
\]
与
\[
\tau \xrightarrow{w} h^{*}(\zeta) \xrightarrow{c_{g,h}(\eta)} (gh)^{*}(\eta) \xrightarrow{(gh)^{*}(u)} (gh)^{*} f^{*}(\xi) \xrightarrow{c_{f,gh}(\xi)} (fgh)^{*}(\xi)
\]
(中间方形交换是因为 $c_{g,h}$ 是函子同态，右边方形按伪函子的条件 B) 交换). 由此得所述结果.

当然，还需在伪函子已来自 $\mathcal{E}$ 上的规范化劈开范畴 $\mathcal{F}'$ 的情形明确如何得到 $\mathcal{F}'$ 与 $\mathcal{F}$ 之间的自然同构. 细节留给读者.

也把劈开范畴 $\mathcal{F}$ (在 $\mathcal{E}$ 上) 经基变换函子 $\mathcal{E}' \to \mathcal{E}$ 的逆像概念用伪函子解释留给读者.

\section{例：由函子 \texorpdfstring{$\mathcal{E}^{\circ} \to \mathbf{Cat}$}{E^o→Cat} 定义的劈开范畴；\texorpdfstring{$\mathcal{E}$}{E} 上的分裂范畴}

假设有函子
\[
\varphi \colon \mathcal{E}^{\circ} \longrightarrow \mathbf{Cat},
\]
它定义伪函子，令
\[
\mathcal{F}(S) = \varphi(S), \quad f^{*} = \varphi(f), \quad c_{f,g} = \mathrm{id}_{(fg)^{*}}.
\]
因此前一节的构造给出 $\mathcal{E}$ 上的劈开范畴 $\mathcal{F}$，称为与函子 $\varphi$ 相关的劈开范畴. 为使 $\mathcal{E}$ 上的劈开范畴同构于由函子 $\varphi \colon \mathcal{E}^{\circ} \to$ Cat 定义的劈开范畴，显然必要且充分条件是它满足条件：
\[
(fg)^{*} = g^{*} f^{*}, \quad c_{f,g} = \mathrm{id}_{(fg)^{*}}.
\]
用运输态射集 $K$ 的术语，这也简单意味着两个运输态射的复合是运输态射. $\mathcal{E}$ 上的范畴 $\mathcal{F}$ 的满足前述条件的劈开称为 $\mathcal{F}$ 在 $\mathcal{E}$ 上的\emph{分裂} (splitting)\index[term]{fenlie@分裂}，配备分裂的 $\mathcal{E}$ 之上的范畴 $\mathcal{F}$ 称为 $\mathcal{E}$ 上的\emph{分裂范畴} (split category)\index[term]{fenliefanchou@分裂范畴}. 因此它是劈开范畴概念的特殊情形. $\mathcal{E}$ 上的分裂范畴范畴于是等价于 Hom$(\mathcal{E}^{\circ}, \mathrm{Cat})$. 注意 $\mathcal{E}$ 上的分裂范畴更是 $\mathcal{E}$ 上的劈开范畴.

若 $\mathcal{F}$ 是 $\mathcal{E}$ 上的纤维范畴，$\mathcal{F}$ 上不一定存在分裂. 例如假设 $\mathrm{Ob} \mathcal{E}$ 和 $\mathrm{Ob} \mathcal{F}$ 约化为一个元素，且该对象的自同态集是群 $E$ (相应地：$F$)，从而投影函子 $p$ 由群同态 $p \colon F \to E$ 给出，由于 $p$ 纤维故满射. 立即验证 $\mathcal{F}$ 在 $\mathcal{E}$ 上的劈开集与映射 $s \colon E \to F$ (满足 $ps = \mathrm{id}_E$，i.e. 满射同态 $p \colon F \to E$ 的核 $G$ 的陪集类的"代表元系") 的集合一一对应. 劈开是分裂当且仅当 $s$ 是群同态. 说存在分裂，意味着 $E$ 以 $G$ 为核的群扩张 $F$ 平凡，当 $G$ 交换时这由 $H^2(E, G)$ 中某个上同调类的消没表达 (其中 $G$ 看作 $E$ 作用其上的群).

然而假设 $\mathcal{F}$ 是 $\mathcal{E}$ 上的纤维范畴，使诸 $\mathcal{F}_S$ 是\emph{刚性范畴} (rigid category)\index[term]{gangxingfanchou@刚性范畴}，i.e. $\mathcal{F}_S$ 的每个对象的自同构群约化为恒等. 则容易证明 $\mathcal{F}$ 在 $\mathcal{E}$ 上有分裂. 事实上，先看到分裂的存在性问题不受把 $\mathcal{F}$ 换成 $\mathcal{E}$-等价范畴的影响，这把我们归结到诸 $\mathcal{F}_S$ 是刚性且约化 (reduced)\index[term]{yuehua@约化} (i.e. $\mathcal{F}_S$ 中两个同构对象相同) 范畴的情形. 但若 $G$ 是刚性约化范畴，则两个函子 $H \to G$ 的任何同构 (其中 $H$ 是任意范畴) 是恒等. 由此得出，若 $\mathcal{F}$ 是 $\mathcal{E}$ 上的纤维范畴，使纤维范畴刚性且约化，则 $\mathcal{F}$ 在 $\mathcal{E}$ 上有唯一劈开，它必然是一个分裂. 故 $\mathcal{F}$ 同构于由函子 $\varphi \colon \mathcal{E}^{\circ} \longrightarrow \mathbf{Cat}$ 定义的范畴，其中诸 $\varphi(S)$ 是刚性且离散的范畴，且函子 $\varphi$ 在唯一同构意义下确定.

\section{余纤维范畴、双纤维范畴}

考虑 $\mathcal{E}$ 之上的范畴 $\mathcal{F}$，投影函子为
\[
p \colon \mathcal{F} \longrightarrow \mathcal{E},
\]
它借助投影函子
\[
p^{\circ} \colon \mathcal{F}^{\circ} \longrightarrow \mathcal{E}^{\circ}
\]
定义 $\mathcal{E}^{\circ}$ 之上的范畴 $\mathcal{F}^{\circ}$.

$\mathcal{F}$ 中的态射 $\alpha \colon \eta \to \xi$ 称为余笛卡儿 (co-cartesian)\index[term]{yudikaertaishe@余笛卡儿态射}，若它是 $\mathcal{E}^{\circ}$ 之上的 $\mathcal{F}^{\circ}$ 的笛卡儿态射. 明确化，看到这意味着：对 $\mathcal{F}_S$ 的每个对象 $\xi'$，映射 $u \mapsto u \circ \alpha$
\[
\mathrm{Hom}_S(\xi, \xi') \longrightarrow \mathrm{Hom}_f(\eta, \xi')
\]
双射. 此时也说 $(\xi, \alpha)$ 是 $\eta$ 经 $f$ 在 $\mathcal{E}$ 之上的范畴 $\mathcal{F}$ 中的正像 (direct image). 若它对每个 $\eta \in \mathcal{F}_T$ 存在，说经 $f$ 的正像函子存在，选定后记为 $f_{*}^{\mathcal{F}}$ 或 $f_{*}$\index[notations]{$f_{*}$}. 它由 $\mathcal{F}_T^{\circ} \times \mathcal{F}_S$ 上的双函子同构定义：
\[
\mathrm{Hom}_S(f_{*}(\eta), \xi) \stackrel{\sim}{\longrightarrow} \mathrm{Hom}_f(\eta, \xi).
\]
故若 $f_{*}$ 存在，为使 $f^{*}$ 存在，必要且充分条件是 $f_{*}$ 有伴随函子，i.e. 存在函子 $f^{*} \colon \mathcal{F}_S \to \mathcal{F}_T$ 和双函子同构
\[
\operatorname{Hom}_S(f_{*}(\eta), \xi) \xrightarrow{\sim} \operatorname{Hom}_T(\eta, f^{*}(\xi)).
\]
设 $g \colon U \to T$ 是 $\mathcal{E}$ 中的另一个态射，假设经 $f$，$g$ 和 $fg$ 的逆像和正像存在. 考虑函子同态
\[
\begin{array}{l}
c^{f,g} \colon f_{*} g_{*} \longleftarrow (fg)_{*} \\
c_{f,g} \colon g^{*} f^{*} \longrightarrow (fg)^{*}.
\end{array}
\]
看到若把 $f_{*} g_{*}$ 和 $g^{*} f^{*}$ 看作一对伴随函子，同样 $(fg)_{*}$ 和 $(fg)^{*}$，则前面两个同态彼此伴随. 故其中一个同构当且仅当另一个同构. 特别：

\begin{proposition}\label{prop:VI.10.1}
假设 $\mathcal{E}$ 之上的范畴 $\mathcal{F}$ 预纤维且余预纤维. 为使它纤维，必要且充分条件是它余纤维.
\end{proposition}

当然，若 $\mathcal{F}^{\circ}$ 在 $\mathcal{E}$ 上预纤维 (相应地：纤维)，则说 $\mathcal{F}$ 在 $\mathcal{E}$ 上余预纤维 (co-prefibered) (相应地：余纤维 (co-fibered)\index[term]{yuxianweifanchou@余纤维范畴}). 若 $\mathcal{F}$ 在 $\mathcal{E}$ 上同时纤维和余纤维，则说 $\mathcal{F}$ 在 $\mathcal{E}$ 上双纤维 (bi-fibered)\index[term]{shuangxianweifanchou@双纤维范畴}.

\section{各种例子}

\begin{enumerate}[a)]
\item \textbf{$\mathcal{E}$ 的箭范畴.} 设 $\mathcal{E}$ 是范畴. 用 $\Delta^{1}$ 记与全序二元集 $[0, 1]$ 相关的范畴；它因此有两个对象 0 和 1，除两个恒等态射外，还有一条从 0 到 1 的箭 $(0, 1)$. 设
\[
\mathbf{Fl}(\mathcal{E}) = \mathbf{Hom}(\Delta^{1}, \mathcal{E}),
\]
称它为 $\mathcal{E}$ 的箭范畴. $\Delta^{1}$ 的对象 1 定义典范函子，称为\emph{靶函子} (target functor)\index[term]{bahanzi@靶函子}
\[
\mathbf{Fl}(\mathcal{E}) \longrightarrow \mathcal{E}
\]
($\Delta^{1}$ 的对象 0 定义的函子称为源函子). 对 $\mathcal{E}$ 的每个对象 $S$，纤维范畴 $\mathbf{Fl}(\mathcal{E})_S$ 典范同构于 $S$ 之上的 $\mathcal{E}$ 的对象范畴 $\mathcal{E}_{/S}$.

考虑 $\mathcal{E}$ 中的态射 $f \colon T \to S$，它对应典范函子
\[
f_{*} \colon \mathcal{E}_{/T} = \mathcal{F}_T \longrightarrow \mathcal{E}_{/S} = \mathcal{F}_S
\]
和函子同构
\[
\operatorname{Hom}_S(f_{*}(\eta), \xi) \xrightarrow{\sim} \operatorname{Hom}_f(\eta, \xi)
\]
它因此使 $f_{*}$ 成为 $\mathcal{F}$ 中经 $f$ 的正像函子. 这里还有
\[
(\mathrm{id}_S)_{*} = \mathrm{id}_{\mathcal{F}_S}, \quad (fg)_{*} = f_{*} g_{*}, \quad c^{f,g} = \mathrm{id}_{(fg)_{*}},
\]
i.e. $\mathcal{F}$ 配备 $\mathcal{E}$ 上的余分裂 (co-splitting). 更不用说，$\mathcal{F}$ 在 $\mathcal{E}$ 上余纤维. 注意 $\mathcal{F}$ 中态射集与 $\mathcal{E}$ 中的交换方形图式一一对应. 按定义，当所述方形在 $\mathcal{E}$ 中笛卡儿，i.e. 使 $Y$ 成为 $X$ 和 $T$ 在 $S$ 上的纤维积时，该态射笛卡儿. 故经 $f$ 的逆像函子 $f^{*}$ 存在，当且仅当对 $S$ 上的每个对象 $X$，纤维积 $X \times_S T$ 存在. 由 \ref{prop:VI.10.1} 得出，若 $\mathcal{E}$ 中第三个之上的两个对象的积总是存在，i.e. 若 $\mathcal{F}$ 在 $\mathcal{E}$ 上预纤维，则 $\mathcal{F}$ 甚至在 $\mathcal{E}$ 上纤维.

\item \textbf{变动空间上的预层或层范畴.} 设 $\mathcal{E} = \mathbf{Top}$ 是拓扑空间范畴. 若 $T$ 是拓扑空间，用 $\mathcal{U}(T)$ 记 $T$ 的开集范畴，其中态射是包含映射. 若 $\mathcal{C}$ 是范畴，从 $\mathcal{U}(T)^{\circ}$ 到 $\mathcal{C}$ 的函子称为 $T$ 上取值于 $\mathcal{C}$ 的\emph{预层} (presheaf)\index[term]{yuceng@预层}，若它满足这里不再重复的左正合条件，则称为层 (sheaf)\index[term]{ceng@层}. $T$ 上取值于 $\mathcal{C}$ 的预层范畴 $\mathcal{P}(T)$ 按定义是范畴 Hom$(\mathcal{U}(T)^{\circ}, \mathcal{C})$，$T$ 上取值于 $\mathcal{C}$ 的层范畴 $\mathcal{F}(T)$ 是 Hom$(\mathcal{U}(T)^{\circ}, \mathcal{C})$ 的由是层的对象组成的全子范畴. 若 $f \colon T \to S$ 是 $\mathcal{E}$ 中的态射，i.e. 拓扑空间的连续映射，由增映射 $U \mapsto f^{-1}(U)$ 对应函子 $\mathcal{U}(S) \to \mathcal{U}(T)$，由此得函子
\[
f_{*} \colon \operatorname{Hom}(\mathcal{U}(T)^{\circ}, \mathcal{C}) \longrightarrow \operatorname{Hom}(\mathcal{U}(S)^{\circ}, \mathcal{C})
\]
称为经 $f$ 的预层正像函子. 立即看到层的正像是层，故函子 $f_{*} \colon \mathcal{P}(T) \to \mathcal{P}(S)$ 诱导函子，仍记为 $f_{*} \colon \mathcal{F}(T) \to \mathcal{F}(S)$. 平凡地验证 (用函子复合的结合性) 对第二个连续映射 $g \colon U \to T$ 有恒等
\[
(gf)_{*} = g_{*} f_{*}, \quad \text{同样 } (\mathrm{id}_S)_{*} = \mathrm{id}_{\mathcal{P}(S)}.
\]
这样，得到函子
\[
S \longmapsto \mathcal{P}(S)
\]
(相应地：
\[
S \longmapsto \mathcal{F}(S)
\])
从 $\mathcal{E}$ 到 Cat. 事实上，我们关心相应函子
\[
S \longmapsto \mathcal{P}(S)^{\circ}, \quad \text{相应地: } S \longmapsto \mathcal{F}(S)^{\circ}.
\]
它定义拓扑空间范畴上的余纤维 (甚至余分裂) 范畴，称为取值于 $\mathcal{C}$ 的预层 (相应地：层) \emph{余纤维范畴} (暗指：在变动空间之上). 明确第 8 号的构造，看到 $T$ 上的预层 $B$ 到 $S$ 上的预层 $A$ 的态射是对 $(f, u)$，由 $T$ 到 $S$ 的连续映射和 $\mathcal{P}(S)$ 中的态射 $u \colon A \to f_{*}(B)$ 组成. 这个描述对层的态射同样有效，$\mathcal{F}$ 是 $\mathcal{P}$ 的全子范畴.

在最重要的情形，$\mathcal{E}$ 之上的范畴 $\mathcal{P}$ 和 $\mathcal{F}$ 也是纤维范畴，i.e. 对每个连续映射，正像函子 $\mathcal{P}(T) \to \mathcal{P}(S)$ 和 $\mathcal{F}(T) \to \mathcal{F}(S)$ 有伴随函子，记为 $f^{*}$，称为经连续映射 $f$ 的预层逆像函子 (相应地：层逆像函子). 例如 $\mathcal{C} = \mathbf{Ens}$ 时该函子存在. 可以证明函子 $f^{*} \colon \mathcal{P}(S) \to \mathcal{P}(T)$ 只要 $\mathcal{C}$ 中存在归纳极限 (相对于所考虑宇宙中的图式) 就存在. 对 $\mathcal{F}$ 的问题较不易；注意 (即使 $\mathcal{C} = \mathbf{Ens}$) 是层的预层的逆像一般不是层，换言之，层逆像函子不同构于预层逆像函子诱导的函子 (尽管共用记号 $f^{*}$). 这样，$\mathcal{F}$ 是 $\mathcal{P}$ 的余纤维子范畴，但不是纤维子范畴，i.e. 包含函子 $\mathcal{F} \to \mathcal{P}$ 不是纤维的.

余纤维范畴 $\mathcal{P}$ 可以由更一般的余纤维 (或更确切地说纤维) 范畴推出，后者如下得到. 对每个范畴 $\mathcal{U}$ (在固定宇宙中)，令
\[
\mathcal{P}(\mathcal{U}) = \operatorname{Hom}(\mathcal{U}, \mathcal{C})
\]
并注意 $\mathcal{U} \mapsto \mathcal{P}(\mathcal{U})$ 自然地是 Cat 到 Cat 的 $U$ 的反变函子. 它因此定义 $\mathcal{E} = \mathbf{Cat}$ 之上的分裂范畴，记为 $\mathbf{Cat}_{//\mathcal{C}}$. 该范畴的对象是范畴 $\mathcal{U}$ 和函子 $p \colon \mathcal{U} \to \mathcal{C}$ 组成的对 $(\mathcal{U}, p)$，$(\mathcal{U}, p)$ 到 $(\mathcal{V}, q)$ 的态射本质上是函子 $f \colon \mathcal{U} \to \mathcal{V}$ 和函子同态 $u \colon p \to qf$ 组成的对 $(f, u)$. 把 $\mathbf{Cat}_{//\mathcal{C}}$ 中态射的复合明确留给读者. 投影函子
\[
\mathcal{F} = \mathrm{Cat}_{//\mathcal{C}} \longrightarrow \mathcal{E} = \mathrm{Cat}
\]
给对 $(\mathcal{U}, p)$ 配对象 $\mathcal{U}$；在 $\mathcal{U}$ 处的纤维范畴是范畴 Hom$(\mathcal{U}, \mathcal{C})$ (在同构意义下). 当 $\mathcal{C}$ 中存在归纳极限时，容易证明 Cat 之上的纤维范畴 $\mathbf{Cat}_{//\mathcal{C}}$ 在 Cat 上也是余纤维的，i.e. 可以定义函子 $p \colon \mathcal{U} \to \mathcal{C}$ 经函子 $f \colon \mathcal{U} \to \mathcal{V}$ 的正像概念. 预层范畴由上述纤维范畴经基变换
\[
\mathbf{Top}^{\circ} \longrightarrow \mathbf{Cat}
\]
(前面定义的函子 $S \mapsto \mathcal{U}(S)$) 推出，得到 $\mathbf{Top}^{\circ}$ 之上的纤维范畴，取对偶范畴得到 Top 之上的预层余纤维范畴 $\mathcal{P}$. 函子的逆像概念对应预层的正像概念，函子的正像概念对应预层的逆像概念.

\item \textbf{算子对象之上的算子对象.} 设 $\mathcal{F}$ 是 $\mathcal{E}$ 之上的范畴，$S$ 是 $\mathcal{E}$ 的对象，有限群 $G$ 在其上 (为确定起见) 左作用. 这个算子对象可以解释为对应函子 $\lambda \colon \mathcal{E}' \to \mathcal{E}$ (其中 $\mathcal{E}'$ 是以 $G$ 为自同态群的单对象范畴) 到 $\mathcal{E}$，并因此由基变换定义 $\mathcal{E}'$ 之上的范畴 $\mathcal{F}'$，当 $\mathcal{F}$ 在 $\mathcal{E}$ 上纤维 (相应地：余纤维) 时 $\mathcal{F}'$ 在 $\mathcal{E}'$ 上纤维 (相应地：余纤维). $\mathcal{E}'$ 在 $\mathcal{F}'$ 上的截面 (必然笛卡儿，因为 $\mathcal{E}'$ 是群胚，且由 \ref{cor:VI.6.12}，$\mathcal{F}'$ 中的每个同构笛卡儿) 也可以解释为 $\lambda$ 之上的 $\mathcal{E}$-函子 $\mathcal{E}' \to \mathcal{F}$，或解释为算子对象 $S_{\circ}$ "之上"的 $\mathcal{F}$ 中的算子对象 $\xi$.

\item \textbf{互为拟逆的伴随函子对；自对偶.} 当基范畴 $\mathcal{E}$ 约化为两个对象 $a$，$b$，且除恒等箭外还有互为逆的两个同构 $f \colon a \to b$ 和 $g \colon b \to a$ (i.e. $\mathcal{E}$ 是带两个对象的连通刚性群胚) 时，$\mathcal{E}$ 上的规范化劈开范畴本质上是两个范畴 $\mathcal{F}_a$ 和 $\mathcal{F}_b$ 及一对伴随函子 $G \colon \mathcal{F}_a \to \mathcal{F}_b$ 和 $F \colon \mathcal{F}_b \to \mathcal{F}_a$ 组成的系统，它们互为范畴等价 (故互为拟逆). 取 $\mathcal{F}_a$ 和 $\mathcal{F}_b$ 为 $\mathcal{F}$ 的纤维范畴，取 $F$ 和 $G$ 为函子 $f^{*}$ 和 $g^{*}$，两个同构
\[
u \colon FG \stackrel{\sim}{\longrightarrow} \mathrm{id}_{\mathcal{F}_a} \qquad v \colon GF \stackrel{\sim}{\longrightarrow} \mathrm{id}_{\mathcal{F}_b}
\]
是 $c_{g,f}$ 和 $c_{f,g}$. $u$ 与 $v$ 之间的两个通常相容性条件正是 $fgf$ 和 $gfg$ 的复合的 \ref{prop:VI.7.4} B) 条件. 容易证明这些条件足以蕴含确实有伪函子 $\mathcal{E}^{\circ} \to \mathbf{Cat}$.

有趣的情形是
\[
\mathcal{F}_b = \mathcal{F}_a^{\circ}, \quad G = F^{\circ}, \quad v = u^{\circ}.
\]
称范畴 $\mathcal{C}$ 中的\emph{自对偶} (self-duality)\index[term]{ziduiou@自对偶} 为函子 $D \colon \mathcal{C} \to \mathcal{C}^{\circ}$ 和同构 $u \colon DD^{\circ} \xrightarrow{\sim} \mathrm{id}_{\mathcal{C}}$ 的给出，使 $u$ 与同构 $u^{\circ} \colon D^{\circ}D \stackrel{\sim}{\longrightarrow} \mathrm{id}_{\mathcal{C}^{\circ}}$ 使 $(D, D^{\circ})$ 成为一对伴随函子 (必然互为拟逆). 该条件写成：
\[
D(u(x)) = u(D(x)) \quad \text{对每个 } x \in \mathrm{Ob}(\mathcal{C}).
\]
\item \textbf{离散范畴 $\mathcal{E}$ 之上的范畴.} 若 $\mathcal{E}$ 的每条箭都是恒等箭，则称 $\mathcal{E}$ 是\emph{离散范畴} (discrete category)\index[term]{lisanfanchou@离散范畴}，此时 $\mathcal{E}$ 在唯一同构意义下由集合 $I = \mathrm{Ob}(\mathcal{E})$ 确定. $\mathcal{E}$ 之上范畴 $\mathcal{F}$ 的给出于是 (在唯一同构意义下) 等价于一族范畴 $\mathcal{F}_i$ $(i \in I)$ 的给出，即纤维范畴. $\mathcal{E}$ 之上的每个范畴 $\mathcal{F}$ 纤维，每个 $\mathcal{E}$-函子 $\mathcal{F} \to \mathcal{G}$ 笛卡儿，有典范同构
\[
\mathbf{Hom}_{\mathcal{E}/-}(\mathcal{F}, \mathcal{G}) \stackrel{\sim}{\longrightarrow} \prod_{i} \mathbf{Hom}(\mathcal{F}_i, \mathcal{G}_i).
\]
特别得到
\[
\Gamma(\mathcal{F}/\mathcal{E}) = \varprojlim \mathcal{F}/\mathcal{E} \stackrel{\sim}{\longrightarrow} \prod_{i} \mathcal{F}_i.
\]
\item 假设 $\mathcal{E}$ 恰有两个对象 $S$ 和 $T$，且除恒等态射外还有态射 $f \colon T \to S$. 则 $\mathcal{E}$ 之上的范畴 $\mathcal{F}$ 在唯一 $\mathcal{E}$-同构意义下由两个范畴 $\mathcal{F}_S$ 和 $\mathcal{F}_T$ 及取值于 Ens 的 $\mathcal{F}_T^{\circ} \times \mathcal{F}_S$ 上的双函子 $H(\eta, \xi)$ 定义. 事实上，若 $\mathcal{F}$ 是 $\mathcal{E}$ 之上的范畴，给它配两个纤维范畴 $\mathcal{F}_S$ 和 $\mathcal{F}_T$ 及双函子 $H(\eta, \xi) = \mathrm{Hom}_f(\eta, \xi)$. 反向构造的明确留给读者. 为使所考虑的范畴纤维 (或预纤维，这在此相同)，必要且充分条件是函子 $H$ 关于自变量 $\xi$ 可表示. 为使它余纤维，必要且充分条件是 $H$ 关于自变量 $\eta$ 可表示.
\item 设 $\mathcal{F} = \mathcal{C} \times \mathcal{E}$，借助 $\mathrm{pr}_2$ 把它看作 $\mathcal{E}$ 之上的范畴. 则 $\mathcal{F}$ 在 $\mathcal{E}$ 上纤维且余纤维，甚至配备由 $\mathcal{E}$ (相应地：$\mathcal{E}^{\circ}$) 上取常值 $\mathcal{C}$ 的函子对应的典范分裂和余分裂. 有
\[
\Gamma(\mathcal{F}/\mathcal{E}) \simeq \operatorname{Hom}(\mathcal{E}, \mathcal{C})
\]
且 Lim $\mathcal{F}/\mathcal{E}$ 对应于由把任意态射变为同构的函子 $F \colon \mathcal{E} \to \mathcal{C}$ 组成的全子范畴.
\end{enumerate}

\section{劈开范畴上的函子}

设 $\mathcal{F}$ 是 $\mathcal{E}$ 上的规范化劈开范畴. 对 $\mathcal{E}$ 的每个对象 $S$，用
\[
i_S \colon \mathcal{F}_S \longrightarrow \mathcal{F}
\]
记包含函子. 对 $\mathcal{E}$ 中的每个态射 $f \colon T \to S$，有函子同态
\[
\alpha_f \colon i_T f^{*} \longrightarrow i_S,
\]
其中 $f^{*}$ 是劈开定义的 $f$ 的基变换函子 $\mathcal{F}_S \to \mathcal{F}_T$. 现在设
\[
F \colon \mathcal{F} \longrightarrow \mathcal{C}
\]
是 $\mathcal{F}$ 到范畴 $\mathcal{C}$ 的函子，对每个 $S \in \mathrm{Ob}(\mathcal{E})$ 令
\[
F_S = F \circ i_S \colon \mathcal{F}_S \longrightarrow \mathcal{C},
\]
对 $\mathcal{E}$ 中的每个 $f \colon T \to S$，
\[
\varphi_f = F * \alpha_f \colon F_T f^{*} \longrightarrow F_S.
\]
这样，对每个函子 $F \colon \mathcal{F} \to \mathcal{C}$，配了一族函子 $F_S \colon \mathcal{F}_S \to \mathcal{C}$ 和一族函子同态 $\varphi_f \colon F_T f^{*} \to F_S$. 这些族满足下列条件：
\begin{enumerate}[a)]
\item $\varphi_{\mathrm{id}_S} = \mathrm{id}_{F_S}$.
\item 对 $\mathcal{E}$ 中的两个态射 $f \colon T \to S$ 和 $g \colon U \to T$，在函子同态方形中有交换性：
\end{enumerate}
\[
\begin{tikzcd}[row sep=large, column sep=large]
F_U g^{*} f^{*} \arrow[r,"{F_U * c_{f,g}}"] \arrow[d,"{\varphi_g * f^{*}}"] & F_U (fg)^{*} \arrow[d,"{\varphi_{fg}}"] \\
F_T f^{*} \arrow[r,"{\varphi_f}"] & F_S.
\end{tikzcd}
\]
第一个关系平凡，第二个关系通过把函子 $F$ 应用于交换图式
\[
\begin{tikzcd}[row sep=large, column sep=large]
g^{*} f^{*}(\xi) \arrow[r,"{c_{f,g}(\xi)}"] \arrow[d,"{\alpha_g(f^{*}(\xi))}"] & (fg)^{*}(\xi) \arrow[d,"{\alpha_{fg}(\xi)}"] \\
f^{*}(\xi) \arrow[r,"{\alpha_f(\xi)}"] & \xi
\end{tikzcd}
\]
(对 $\mathcal{F}_S$ 中变动的对象 $\xi$) 得到.

若 $G$ 是第二个函子 $\mathcal{F} \to \mathcal{C}$，产生函子 $G_S \colon \mathcal{F}_S \to \mathcal{C}$ 和函子同态 $\psi_f \colon G_T f \to G_S$，且 $u \colon F \to G$ 是函子同态，则它对应函子同态 $u * i_S$：
\[
u_S \colon F_S \longrightarrow G_S,
\]
并立即看到对 $\mathcal{E}$ 中的每个态射 $f \colon T \to S$，在方形中交换
\begin{equation}
\begin{tikzcd}[row sep=large, column sep=large]
F_T f^{*} \arrow[r,"{\phi_f}"] \arrow[d,"{u_T * f^{*}}"] & F_S \arrow[d,"{u_S}"] \\
G_T f^{*} \arrow[r,"{\psi_f}"] & G_S
\end{tikzcd} \tag{c}
\label{eq:VI.c}
\end{equation}

\begin{proposition}\label{prop:VI.12.1}
设 $\mathcal{H}(\mathcal{F}, \mathcal{C})$ 是对象为"函子 $F_S \colon \mathcal{F}_S \to \mathcal{C}$ 的族 $(F_S)$ ($S \in \mathrm{Ob}(\mathcal{E})$) 与函子同态 $F_T f^{*} \to F_S$ 的族 $(\varphi_f)$ ($f \in \mathrm{Fl}(\mathcal{E})$) 组成的对、满足条件 a) 和 b)"、态射为满足前面写出的交换条件 c) 的同态族 $(u_S)$ ($S \in \mathrm{Ob}(\mathcal{E})$) 的范畴 (态射的复合通过函子同态 $\mathcal{F}_S \to \mathcal{C}$ 的复合进行). 则上面明确的两个法则定义 Hom$(\mathcal{F}, \mathcal{C})$ 与范畴 $\mathcal{H}(\mathcal{F}, \mathcal{C})$ 的同构 $K$.
\end{proposition}

显然得到第一个范畴到第二个的函子. 该函子全忠实，因为对给定的 $F$，$G$，映射
\[
\mathrm{Hom}(F, G) \longrightarrow \mathrm{Hom}(K(F), K(G))
\]
平凡单射；为证满射，只需注意交换条件 c) 表达映射 $u(\xi) = u_S(\xi) \colon F(S) = F_S(\xi) \to G(\xi) = G_S(\xi)$ 对 $\mathcal{F}$ 中形如 $\alpha_f(\xi)$ 的同态的函子性，另一方面在每个纤维范畴上 (i.e. 对 $\mathcal{F}$ 中是 $T$-态射 ($T \in \mathrm{Ob}(\mathcal{E})$) 的态射) 有函子性，故对 $\mathcal{F}$ 中的每个态射有函子性，因为 $f$-态射 ($f \colon T \to S$ 是 $\mathcal{E}$ 中的态射) 唯一地是 $\alpha_f(\xi)$ 与 $T$-态射的复合. 剩下证明函子 $K$ 对对象双射. 前面的论证已表明 $K$ 对对象单射，剩下证明满射，i.e. 若从满足 a) 和 b) 的系统 $(F_S), (\varphi_f)$ 出发，用
\[
F(\xi) = F_S(\xi) \quad \text{对 } \xi \in \mathrm{Ob} \mathcal{F}_S \subset \mathrm{Ob} \mathcal{F}
\]
定义映射 $\mathrm{Ob} \mathcal{F} \to \mathrm{Ob} \mathcal{C}$，用
\[
F(\alpha_f(\xi) u') = \varphi_f(\xi) F_T(u')
\]
(对 $\mathcal{E}$ 中的每个态射 $f \colon T \to S$，$\mathcal{F}_S$ 的每个对象 $\xi$ 和以 $f^{*}(\xi)$ 为靶的每个 $T$-态射 $u'$) 定义映射 $\mathrm{Fl}(\mathcal{F}) \to \mathrm{Fl}(\mathcal{C})$，则得到 $\mathcal{F}$ 到 $\mathcal{C}$ 的函子 $F$. 事实上，关系 $F(\mathrm{id}_{\xi}) = \mathrm{id}_{F(\xi)}$ 平凡，剩下在有一个 $f$-态射 $u \colon \eta \to \xi$ 和一个 $g$-态射 $v \colon \zeta \to \nu$ (其中 $f \colon T \to S$，$g \colon U \to T$ 是 $\mathcal{E}$ 的态射) 时证明乘性 $F(uv) = F(u)F(v)$. 令 $w = uv$，有
\[
u = \alpha_f(\xi) u', v = \alpha_g(\eta) v', w = \alpha_{fg}(\xi) w'
\]
其中
\[
w' = c_{f,g}(\xi) g^{*}(u') v' \quad (\text{参见第 8 号}).
\]
用这些记号，要证明下面图式外轮廓的交换性：
\[
\begin{tikzcd}[row sep=large, column sep=large]
F_U(\zeta) \arrow[rrr, phantom, "F(w')"] \arrow[r,"{F_U(v')}"] \arrow[d,"{\varphi_g(\eta)}"] & F_U g^{*}(\eta) \arrow[r,"{F_U g^{*}(u')}"] \arrow[d,"{\varphi_g(f^{*}(\xi))}"] & F_U g^{*} f^{*}(\xi) \arrow[r,"{F_U(c_{f,g}(\xi))}"] \arrow[d,"{\varphi_{fg}(\xi)}"] & F_U (fg)^{*}(\xi) \\
F_T(\eta) \arrow[rr, phantom, "F(u)"] \arrow[r,"{F_T(u')}"] & F_T f^{*}(\xi) \arrow[r,"{\varphi_f(\xi)}"] & F_S(\xi)
\end{tikzcd}
\]
左边三角形按 $F(v)$ 的定义交换，中间方形按 $u' \colon \xi \to f(\eta)$ 经函子同态 $\alpha_g$ 得出的定义交换，右边方形按条件 b) 交换. 结论由此得出.

现在假设 $\mathcal{C}$ 也是 $\mathcal{E}$ 上的规范化劈开范畴，此后称它为 $\mathcal{G}$，且关心 $\mathcal{F}$ 到 $\mathcal{G}$ 的 $\mathcal{E}$-函子. 若 $F$ 是这样一个函子，它对纤维范畴诱导函子
\[
F_S \colon \mathcal{F}_S \longrightarrow \mathcal{G}_S.
\]
另一方面，对 $\mathcal{E}$ 中的每个态射 $f \colon T \to S$ 和 $\mathcal{F}_S$ 的每个对象 $\xi$，$f$-态射 $F(\alpha_f(\xi))$ 唯一地经 $T$-态射
\[
\varphi_f(\xi) \colon F_T(f_{\mathcal{F}}^{*}(\xi)) \longrightarrow f_{\mathcal{G}}^{*}(F_S(\xi))
\]
分解 (其中下标 $\mathcal{F}$ 或 $\mathcal{G}$ 指明取逆像函子的劈开范畴)，由此得 $\mathcal{F}_S$ 到 $\mathcal{G}_T$ 的函子同态：
\[
\varphi_f \colon F_T f_{\mathcal{F}}^{*} \longrightarrow f_{\mathcal{G}}^{*} F_S.
\]
两个系统 $(F_S)$ 和 $(\varphi_f)$ 满足下列条件：
a') $\varphi_{\mathrm{id}_S} = \mathrm{id}_{F_S}$.
b') 对 $\mathcal{E}$ 中的两个态射 $f \colon T \to S$ 和 $g \colon U \to T$，在下列函子同态图式中有交换性：
\[
\begin{tikzcd}[row sep=large, column sep=large]
F_U g_{\mathcal{F}}^{*} f_{\mathcal{F}}^{*} \arrow[r,"{F_U * c^{\mathcal{F}}_{f,g}}"] \arrow[d,"{\varphi_g f_{\mathcal{F}}^{*}}"] & F_U (fg)_{\mathcal{F}}^{*} \\
g_{\mathcal{G}}^{*} F_T f_{\mathcal{F}}^{*} \arrow[d,"{g_{\mathcal{G}}^{*} * \varphi_f}"] & \\
g_{\mathcal{G}}^{*} f_{\mathcal{G}}^{*} F_S \arrow[r,"{c^{\mathcal{G}}_{f,g} * F_S}"] & (fg)_{\mathcal{G}}^{*} F_S.
\end{tikzcd}
\]
验证留给读者，\cref{prop:VI.12.1} 的类似陈述的表述和证明也留给读者，它蕴含 $\mathcal{F}$ 到 $\mathcal{G}$ 的 $\mathcal{E}$-函子集与满足上面条件 a') 和 b') 的系统 $(F_S), (\varphi_f)$ 集之间的一一对应. 当然，在这个对应中，笛卡儿函子由"诸同态 $\varphi_f$ 是同构"这一性质刻画.

\begin{remarkstar}
当然，大多数时候直接在纤维范畴上推理而不使用显式劈开是有利的，这尤其使简单概念 $\mathcal{E}$-函子或笛卡儿 $\mathcal{E}$-函子不需要上面那样笨重的解释. 正是为了避免难以忍受的笨拙、并得到更内蕴的陈述，我们不得不放弃像 \cite{exp06-2} 那样从劈开范畴 (loc. cit. 中称为"纤维范畴") 的概念出发，后者退居其次，让位于纤维范畴概念. 而且很可能，与目前仍占优势、与旧思维习惯相关的用法相反，在万有问题中最终将被证明更便利的，不是把重点放在某个假定一次选定好的解上，而是让所有解处于平等地位.
\end{remarkstar}

\section{文献}\printbibliography[heading=none]


